% BEGIN EMBEDDED BIBLIOGRAPHY
% Bibliography data are part of this source; the compiler writes its temporary input.
\begin{filecontents*}[overwrite]{\jobname.bib}
@unpublished{QED3Master,
  title = {The field-strength {Osterwalder--Schrader} axioms for {QED} in three dimensions},
  date = {2026-10-20},
  note = {Comprehensive edition, 263 pp., 20 October 2026, before the modular integration revision}
}
@article{DimockUVS2022,
  author = {Dimock, J.},
  title = {Ultraviolet stability for {QED} in {$d=3$}},
  journal = {Annales Henri Poincar\'e},
  volume = {23},
  date = {2022},
  pages = {2113--2205},
  doi = {10.1007/s00023-021-01127-z},
  eprint = {2009.01156},
  eprinttype = {arxiv}
}
@online{PongeSchatten2025,
  author = {Ponge, Rapha\"el},
  title = {Singular Values and Schatten Classes},
  date = {2025},
  note = {Introduction to Noncommutative Geometry, Chapter 4; slides 3--9 and 53--56},
  url = {https://raphaelponge.org/wp-content/uploads/2025/07/summer-chap4-1.pdf}
}

@unpublished{ChamberLattice,
 title={Reflection positivity for chamber {QED} on the infinite lattice},
 date={2026-10-20}, note={Companion manuscript, 20 October 2026}}

@unpublished{CoverTuning,
 title={Agreement across covers and common tuning for lattice {QED} in three dimensions},
 date={2026-10-20}, note={Companion manuscript, 20 October 2026}}

@unpublished{FermionBlock,
 title={Chamber fermion locality and local integration},
 date={2026-10-20}, note={Companion manuscript, 20 October 2026}}

@unpublished{MaxwellBlock,
 title={Maxwell block minimization and Gaussian fluctuations},
 date={2026-10-20}, note={Companion manuscript, 20 October 2026}}

@unpublished{NonlinearStep,
 title={The exact nonlinear renormalization step for lattice {QED} in three dimensions},
 date={2026-10-20}, note={Companion manuscript, 20 October 2026}}

@unpublished{NormalizedContraction,
 title={Relevant extraction and contraction of local interactions},
 date={2026-10-20}, note={Companion manuscript, 20 October 2026}}

@unpublished{QED3Integration,
 title={The continuum field-strength construction for {QED} in three dimensions},
 date={2026-10-20}, note={Modular revision of 20 October 2026}}

@unpublished{ReferenceResponse,
 title={Exact source identities and normalized reference response for lattice {QED} in three dimensions},
 date={2026-10-20}, note={Companion manuscript, 20 October 2026}}

@unpublished{PositiveProbe,
 title={Positive compact probes and free response coefficients for lattice {QED} in three dimensions},
 date={2026-10-20}, note={Companion manuscript, 20 October 2026}}
\end{filecontents*}
% END EMBEDDED BIBLIOGRAPHY
\documentclass[11pt,letterpaper,reqno]{amsart}
% BEGIN INLINE manuscript-style
\makeatletter
\NeedsTeXFormat{LaTeX2e}[2020/10/01]
% Use with \documentclass[11pt,letterpaper,reqno]{amsart}.
% This file controls presentation only. It supplies no mathematical content.

\RequirePackage[T1]{fontenc}
\RequirePackage{lmodern}
\RequirePackage{microtype}
\RequirePackage[letterpaper,margin=1in,headheight=15pt]{geometry}
\RequirePackage{amsmath,amsthm,amssymb,mathtools}
\RequirePackage{enumitem}
\RequirePackage{perpage}
\RequirePackage{xcolor}
\definecolor{manuscriptred}{rgb}{0.6,0.1,0.1}
\definecolor{manuscriptblue}{rgb}{0.1,0.3,0.5}
\definecolor{manuscriptmagenta}{HTML}{60445C}
\RequirePackage[breaklinks=true,colorlinks=true,
  linkcolor=manuscriptred,citecolor=manuscriptblue,urlcolor=manuscriptmagenta]{hyperref}
\RequirePackage{cleveref}
\hypersetup{pdfauthor={}}

% Retain the AMS paragraph, title, heading and running-head conventions.
% Equations retain one global Arabic sequence, as in the reference.
\setcounter{tocdepth}{1}
\renewcommand{\thefigure}{\Roman{figure}}
\renewcommand{\thetable}{\Roman{table}}
\MakePerPage{footnote}
\renewcommand{\thefootnote}{%
  \ifnum\value{footnote}>9\arabic{footnote}\else\fnsymbol{footnote}\fi}

% Upright statements. Main results are A, B, C; supporting statements
% share a section-local sequence. These end marks indicate typography,
% never proof completeness, validation, or a deferred-proof status.
\theoremstyle{definition}
\newtheorem{ms@theorem}{Theorem}
\renewcommand{\thems@theorem}{\Alph{ms@theorem}}
\newtheorem*{ms@theoremstar}{Theorem}
\newtheorem{ms@proposition}{Proposition}[section]
\newtheorem{ms@lemma}[ms@proposition]{Lemma}
\newtheorem{ms@definition}[ms@proposition]{Definition}
\newtheorem{ms@corollary}[ms@proposition]{Corollary}
\newtheorem*{ms@corollarystar}{Corollary}

\newcommand{\ms@majorstart}{%
  \pushQED{\qed}\renewcommand{\qedsymbol}{$\blacksquare\square$}}
\newcommand{\ms@minorstart}{%
  \pushQED{\qed}\renewcommand{\qedsymbol}{$\blacksquare$}}
\newenvironment{theorem}[1][]{\ms@majorstart\ms@theorem[#1]}
  {\popQED\endms@theorem}
\newenvironment{theorem*}[1][]{\ms@majorstart\ms@theoremstar[#1]}
  {\popQED\endms@theoremstar}
\newenvironment{proposition}[1][]{\ms@minorstart\ms@proposition[#1]}
  {\popQED\endms@proposition}
\newenvironment{lemma}[1][]{\ms@minorstart\ms@lemma[#1]}
  {\popQED\endms@lemma}
\newenvironment{definition}[1][]{\ms@minorstart\ms@definition[#1]}
  {\popQED\endms@definition}
\newenvironment{corollary}[1][]{\ms@majorstart\ms@corollary[#1]}
  {\popQED\endms@corollary}
\newenvironment{corollary*}[1][]{\ms@majorstart\ms@corollarystar[#1]}
  {\popQED\endms@corollarystar}
% The source's lemma/proposition variants are deliberately not carried over:
% choose one marker for each statement type rather than encode proof status.

\theoremstyle{remark}
\newtheorem{remark}[ms@proposition]{Remark}
\newtheorem*{remark*}{Remark}
% The ordinary amsthm proof environment retains its open square.

\crefname{ms@theorem}{theorem}{theorems}
\Crefname{ms@theorem}{Theorem}{Theorems}
\crefname{ms@proposition}{proposition}{propositions}
\Crefname{ms@proposition}{Proposition}{Propositions}
\crefname{ms@lemma}{lemma}{lemmas}
\Crefname{ms@lemma}{Lemma}{Lemmas}
\crefname{ms@definition}{definition}{definitions}
\Crefname{ms@definition}{Definition}{Definitions}
\crefname{ms@corollary}{corollary}{corollaries}
\Crefname{ms@corollary}{Corollary}{Corollaries}
\crefname{remark}{remark}{remarks}
\Crefname{remark}{Remark}{Remarks}

% Bibliography configuration for this manuscript.
  \RequirePackage[backend=biber,style=alphabetic,backref=true,
    giveninits=true,isbn=false,date=year]{biblatex}
  \addbibresource{\jobname.bib}
  \setcounter{biburlnumpenalty}{100}
  \setcounter{biburllcpenalty}{100}
  \setcounter{biburlucpenalty}{100}
  \renewcommand*{\bibfont}{\footnotesize}
  \renewbibmacro{in:}{}
  \DeclareFieldFormat{pages}{#1}
  % Use the public alphabetic style. No private author-hash overrides;
  % repeated author names are printed in full.
\makeatother
% END INLINE manuscript-style
\usepackage{mathrsfs}
\newcommand{\R}{\mathbb R}
\newcommand{\C}{\mathbb C}
\newcommand{\Z}{\mathbb Z}
\newcommand{\N}{\mathbb N}
\newcommand{\dd}{\mathrm d}
\DeclareMathOperator{\supp}{supp}
\DeclareMathOperator{\tr}{tr}
\DeclareMathOperator{\diag}{diag}
\allowdisplaybreaks
\hypersetup{pdftitle={The terminal measure and pressure estimates for lattice QED in three dimensions},pdfauthor={}}
\title[Terminal measure and pressure for lattice QED]{The terminal measure and pressure estimates\\
for lattice QED in three dimensions}
\date{}
\begin{document}
\begin{abstract}
We identify the terminal measure in the block construction of
three-dimensional lattice QED with chamber fermions and prove
a common-pressure estimate on its prescribed family of tori.
The measure identity holds for arbitrary bounded Borel
observables and retains the normalization of every restricted
Gaussian integral. A zero-flux Gaussian bound controls the
large-field complement. A mesh-uniform determinant estimate
then yields a partition comparison for two reference tori,
with a bound independent of the ultraviolet cutoff.
\end{abstract}
\maketitle
\tableofcontents

\section{Introduction}

We use the chamber regulator with weights $(6,6\sqrt2,6\sqrt2,6)$
and Wilson parameter $0<\rho<1$. On the unit cubic torus $C_n$
with $n$ sites per side, let
$\mathcal F_n=dC^1(C_n;\R)$ be the space of exact real two-cochains,
with counting inner product and standard Gaussian measure $\gamma_n$.
Write $A(u)$ for the mean-zero coclosed potential of $u\in\mathcal F_n$.
For $q>0$ and a real lattice mass $\boldsymbol\mu$, define the matter factor
\begin{equation}\label{eq:tpm-matter-factor}
 z_{r,n}(q,\boldsymbol\mu)=
 \int_{[-\pi,\pi]^3}\frac{\dd\theta}{(2\pi)^3}
 \int_{\mathcal F_n}
 \det\bigl(D^{\rm ch}[q(A(u)+H_\theta)]+\boldsymbol\mu\bigr)
 \,\dd\gamma_n(u),\qquad
 H_{\theta,\nu}=\frac{\theta_\nu}{qn}.
\end{equation}
Here $D^{\rm ch}[qB]$ is the unit-mesh, four-spin-component
chamber Dirac operator of \cite[\S2]{ChamberLattice}, with $48$
chamber labels per site and antiperiodic spinors. Its argument $qB$
is the bond phase; the subscript $r$ in $z_{r,n}$ denotes the regulator.
The harmonic average in \eqref{eq:tpm-matter-factor} is the
zero-flux matter factor of \cite[(37)--(40)]{ReferenceResponse}.

Fix once and for all a choice of the block construction constants
satisfying \cite[Proposition~4.1]{QED3Integration}, through its
charge ceilings $e_1\le e_0<1$ and terminal scalar conditions.
Thus $L$ is odd, $M$ is a power of $L$, $M\ge L^2$,
$0<\varepsilon\le1/200$, and the exponents are
$r=2$, $p_0=9$, $p=19$.
The fixed ceilings denoted by $e_0,e_1$ here are the
$\bar e_0,\bar e_1$ of \cite{QED3Integration};
they are distinct from the running charge $e_j$ at $j=0$.
This choice includes the photon and fermion primitive bounds,
the exact normalized block step, and its field domains.
For $0<\bar m\le1$ and $N\in\Z_{\ge0}$ put
\[
 h=L^{\lfloor\log_L\bar m^{-1}\rfloor},\qquad
 \eta=L^{-N},\qquad K=N+\log_Lh,\qquad
 a_j=L^{j-N},\quad e_j=e\sqrt{a_j},\quad \mu_j=\bar m a_j.
\]
We distinguish the initial charge $q=e\sqrt\eta$ from the fixed
ceiling $e_0$. Set $v_j=-\log e_j$, $p_j=v_j^{19}$,
$p_{0,j}=v_j^9$, and let $r_j$ be the least nonnegative
integer power of $L$ which is at least $v_j^2$. Define
\[
 R_{\rm top}=hMr_K,\qquad R_*=R_{\rm top}/L,\qquad
 \mathcal T(e)=\{hL^t:t\in\Z_{\ge0},\ hL^t\le R_{\rm top}\}.
\]
For $R\in\mathcal T(e)$, the threshold $N_{\min}(R,e)$
and first stopping level $k_*(R)\ge1$ are those of
\cite[Proposition~6.1 and (52)]{CoverTuning}:
the stop is the first level with $\ell_j=Mr_j$, where
$\ell_j=R/a_j$. At cutoffs $N\ge N_{\min}(R,e)$, write
$a_*=a_{k_*}$, $e_*=e_{k_*}$, $\mu_*=\mu_{k_*}$ and
$\ell_*=\ell_{k_*}$. These physical quantities are independent of $N$,
and
\begin{equation}\label{eq:tpm-stop-scales}
 Mv_*^2\le\ell_*<LMv_*^2,\qquad
 \mu_*^{-1}<L\ell_*,\qquad 0<e_*\le e_K=e\sqrt h,\qquad
 0<\mu_*\le1.
\end{equation}
The threshold is finite and nonincreasing in $R$.
The numerical consequences of the fixed scheme used below are
listed and verified in \cref{subsec:term-remainder}.

Use the real pair
$P_N(e)=(m_{0,N}(e),\varepsilon_{0,N}^{\rm v}(e))$
of \cite[Theorem~D]{CoverTuning}: it is selected on
$R_{\rm top}$ with zero initial activity and zero terminal
mass shift and vacuum term there. The same pair gives a controlled
trajectory on every $R\in\mathcal T(e)$ above its stated threshold.
In particular,
$|m_{0,N}(e)|\le\frac18 q^{3/4-8\varepsilon}$.
Set $\boldsymbol\mu_N=\bar m\eta+m_{0,N}(e)$ and abbreviate
$z_{r,n}=z_{r,n}(q,\boldsymbol\mu_N)$, $n=R/\eta$.
The total mass $\boldsymbol\mu_N$ may have either sign.
Its free comparison uses the positive reference mass:
\[
 z_{r,n}^{\rm free}
   =\int_{[-\pi,\pi]^3}
        \det(D^{\rm ch}[qH_\theta]+\bar m\eta)
                      \frac{\dd\theta}{(2\pi)^3}.
\]
Let $\mathcal D_m(p)$ denote the $192\times192$ chamber Fourier
symbol at mass $m>0$, and set
\begin{equation}\label{eq:term-common-pressure}
 p_{\rm ch}^0(m)=\int_{[-\pi,\pi]^3}
                   \log\det\mathcal D_m(p)\frac{\dd p}{(2\pi)^3},
 \qquad
 p_F=p_{\rm ch}^0(\bar m\eta)+\varepsilon_{0,N}^{\rm v}(e).
\end{equation}
The positive constants $C_6,\kappa_{(0)}$, depending only on $\rho$
and the fixed weights, are given in
\cref{prop:term-free-pressure}. Write
$\varepsilon_5(x)=C_6\log(1-e^{-\kappa_{(0)}x})^{-1}$ and
$N_{\rm ref}(e)=N_{\min}(R_*,e)$.
\begin{theorem}[Common pressure and the two-torus comparison]\label{thm:term-common-partition}
Fix the construction and the common pair specified above, and
assume $0<e\sqrt h\le e_1$. For every $R\in\mathcal T(e)$ and
$N\ge N_{\min}(R,e)$, the matter factors $z_{r,n}$ and $z_{r,n}^{\rm free}$ are positive, and
\begin{align}
 \left|\log z_{r,n}-\log z_{r,n}^{\rm free}
                      -\varepsilon_{0,N}^{\rm v}(e)n^3\right|
          &\le3e_*^{1/4-9\varepsilon},\nonumber\\
 \left|\log z_{r,n}-p_Fn^3\right|
          &\le\varepsilon_P(R)
             :=\varepsilon_5(\bar mR)+3e_*^{1/4-9\varepsilon}
          \le\varepsilon_5(1/L)+3e_K^{1/4-9\varepsilon}.
                                     \label{eq:term-common-pressure-bound}
\end{align}
In particular let $R_*=R_{\mathrm{top}}/L$,
$n_*=R_*/\eta$, and $N\ge N_{\rm ref}(e)$.
The two tori $R_*$ and $LR_*$ use the same mass and
the same pressure. Their matter comparison satisfies
\begin{align}
 |D_{\rm mat}|
 &:=\left|L^3\log z_{r,n_*}-\log z_{r,Ln_*}\right|
       \le L^3\varepsilon_P(R_*)+\varepsilon_P(LR_*)
       \nonumber\\
 &\le (L^3+1)
          \left[\varepsilon_5(\bar mR_*)
                                      +3e_K^{1/4-9\varepsilon}\right]
       \le C_{\rm mat},\nonumber\\
 C_{\rm mat}&=(L^3+1)\bigl(\varepsilon_5(1/L)+1\bigr).
                                     \label{eq:term-matter-comparison}
\end{align}
The same $p_F=p_F(N,e,\bar m)$ is used for every asserted side.
The displayed error bounds are uniform in $N$; $C_{\rm mat}$
depends only on $L,\rho$ and the fixed chamber weights.
\end{theorem}

The theorem compares volumes within one fixed construction.
Its pressure $p_F$ may depend on the cutoff. The estimate
does not require a limit of that pressure: the common bulk term
cancels exactly between the two reference tori.

The proof is the terminal integration argument for the block
construction in \cite[\S21--22]{QED3Master}. Its block averaging
and background-field methods belong to the approach of
Ba{\l}aban and collaborators; Dimock develops this approach for
QED in three dimensions in \cite[\S1--2]{DimockUVS2022}.
The chamber determinant, harmonic averaging, and comparison
proved here are those of \cite[\S22]{QED3Master}.
We use the previously proved actual block transformation and
common tuning in \cite{MaxwellBlock,FermionBlock,NormalizedContraction,NonlinearStep,CoverTuning}.

Two estimates connect that transformation to the finite-volume
matter factor. First, Gaussian domination in the zero-flux
sector bounds the probability of leaving any fluctuation box
or the terminal small-curvature set. The sector projection
remains inside the trace proof. Second, the exact change from
curvature to all photon fluctuation coordinates identifies the
restricted measure with a terminal fermion integral. The Gaussian
normalization and the restriction to each box are kept in this
identity, which holds for every bounded Borel observable.

On the terminal small-field set the interacting integral differs
from a positive twisted free determinant by a relative error
of order $e_*^{1/4-9\varepsilon}$. The coefficient-norm estimates
which prove this are collected in
\cref{subsec:interacting-terminal-operators}.
The free determinant has an explicit bulk pressure; a separate
Bloch estimate makes the twist dependence of the block penalty
exponentially small in $\ell_*$, uniformly in its fine mesh.
Combining these bounds with the positive small-field mass gives
two-sided partition estimates. Taking their logarithms and
bounding the terminal vacuum proves \cref{thm:term-common-partition};
the closing argument is in \cref{subsec:tpm-common-proof}.

The generic differentiated block estimates are recorded in
\cref{sec:tpm-marked} for their use in terminal response
estimates. The fixed reference-mass interval used there is
separate from the variable range $0<\mu_*\le1$ in the pressure proof.
\section{Gaussian domination in the zero-flux sector}
\label{subsec:interacting-zero-flux-Gaussian}

The complement estimate below concerns the normalized
zero-flux law, denoted $\bar\mu$: on a cubic torus it is the
probability measure obtained by normalizing the integrand of
\eqref{eq:tpm-matter-factor} after harmonic averaging.
Write $F^{\rm u}$ for the random curvature in unit-lattice
coordinates and $F_N=dC^1(C(n_1,n_2,n_3);\R)$ for its exact
curvature space; on the cubic torus this is $\mathcal F_n$.
The all-flux Gaussian bound does not directly
give a bound for this conditional law.
We therefore retain the sector projection in
the coordinate-slice trace.

Let $Z^{(0)}[\Phi]$ be the unnormalized zero-flux partition
functional with curvature observable $\Phi$, as in
\cite[(34)]{ReferenceResponse}. Put $r_q=2\pi/q$ and
$g_1(t)=(2\pi)^{-1/2}e^{-t^2/2}$.
On $C(n_1,n_2,n_3)$ use the spaces $H^s,H^l$,
$P_G$, $E$, and $M_0$ of
\cite[(80)]{ChamberLattice} and
\cite[(82)]{ChamberLattice}.
Here $H^s,H^l$ are the slice Hilbert spaces at site and link
planes, $P_G$ is the gauge-invariant projection, $E$ is the
electric momentum operator, and $M_0$ is the half-slab
multiplier including the magnetic Gaussian factor.
Let $P_k$ project onto the magnetic-flux
summand $k$. Choose coordinate slice cycles
$S_2,S_3$, through a fixed slice site, and
for $\xi\in[0,1]^2$ set
$b_\xi=q(\xi_2\mathbf1_{S_2}+\xi_3\mathbf1_{S_3})$.
For bounded magnetic factors $u_j$ and
electric factors $v_j$ with sub-Gaussian growth,
let $K_j^\xi$ have the periodized convolution kernel
\[
 v_j(A'-A)e^{i\langle b_\xi,A'-A\rangle}
                         \prod_bg_1(A'(b)-A(b)).
\]
The zero-flux trace formula is
\begin{equation}\label{eq:interacting-zero-flux-Bloch}
 Z^{(0)}\left[\prod_ju_j(F_j^{\rm mag})v_j(F_j^{\rm el})\right]
 =(2\pi)^{n_1n_2n_3}\int_{[0,1]^2}
 \tr\left(P_0\prod_{j=n_1-1,\ldots,0}
                  r_q^{n_2n_3}K_j^\xi P_GM_0u_j(F^{\rm mag})M_0^*\right)
                                                        \,\dd\xi .
\end{equation}
To derive it, sum the electric curvature along
$S_2,S_3$ and over time. The sums are
$r_qk_{12}$ and $r_qk_{13}$.
The added electric phases are therefore
$e^{2\pi i(\xi_2k_{12}+\xi_3k_{13})}$.
Restricting the first magnetic slice to
flux $k_{23}$ inserts $P_{k_{23}}$ in the
word trace, since the sum of $d_sA$ on
the slice vanishes. All operators preserve
that summand. The flux sum is absolutely
convergent by the curvature Gaussian bound
on the modulus of the integrand.
Set $k_{23}=0$ and integrate over $\xi$;
Fourier orthogonality keeps precisely
$k_{12}=k_{13}=0$. This proves
\eqref{eq:interacting-zero-flux-Bloch}.

\begin{proposition}\label{prop:interacting-zero-flux-domination}
For the harmonic-average zero-flux law at every
real lattice mass, every coordinate side
$n_\mu\ge4$, and every real two-cochain $g$,
\begin{equation}\label{eq:interacting-zero-flux-Gaussian}
 \bar\mu(e^{\langle g,F^{\rm u}\rangle})
                  \le e^{\|\operatorname{Pr}_{F_N}g\|^2/2}.
\end{equation}
In particular, for $\phi\in F_N$ with
$\|\phi\|\le1$ and $a\ge0$,
$\bar\mu\{|\langle\phi,F^{\rm u}\rangle|>a\}\le2e^{-a^2/2}$.
\end{proposition}

\begin{proof}
For a real two-cochain $\beta$, put $\mathcal Z(\beta)=
Z^{(0)}[e^{-\langle\beta,F\rangle-\|\beta\|^2/2}]$.
It is finite by the all-flux bound and is
positive because the sector contains the
$z=0$ term, which has positive integral;
the shifted photon weight is strictly positive.
For time direction $1$, split $\beta$
into electric parts $\beta_j^{\rm el}$ and
magnetic parts $\beta_j^{\rm mag}$.
Take the corresponding exponential $u_j$
and $v_j$ in
\eqref{eq:interacting-zero-flux-Bloch}, replacing
$u_j$ first by $u_j^\delta=\min(u_j,\delta^{-1})$.
The bounded magnetic factors make the
word trace applicable.

Completing the electric Gaussian gives
\[
 K_j^\xi=z_jU_j e^{-|E-b_\xi|^2/2},\qquad
 U_j=e^{i\langle\beta_j^{\rm el},E\rangle},\qquad
 |z_j|=1 .
\]
Let $K_\xi=r_q^{n_2n_3}e^{-|E-b_\xi|^2/2}P_G$.
It is positive, and its square root is
$r_q^{n_2n_3/2}e^{-|E-b_\xi|^2/4}P_G$.
The operator $K_\xi^{1/2}M_0$ is Hilbert--Schmidt.
The periodized square-root Gaussian has the
additional phase $e^{i\langle b_\xi,A'-A\rangle}$;
its modulus is bounded by the kernel majorant
used in \cite[(83)]{ChamberLattice}.
The same real-slice integral and flux sum
therefore prove this assertion uniformly
in $\xi$. Multiplication by bounded
$(u_j^\delta)^{1/2}$ preserves it.

Put $Y_j=K_\xi^{1/2}M_0(u_j^\delta)^{1/2}$.
Cyclicity turns the trace into
$\prod_jz_j$ times the closed chain
\[
 \tr\prod_j U_jP_0Y_j(P_0Y_j)^*.
\]
All factors preserve the magnetic-flux
summand; $U_j$ is unitary. Schatten Hölder
\cite[slides~53--56]{PongeSchatten2025}
with $2n_1$ factors of exponent $2n_1$
bounds this trace by
\[
 \prod_j\left(\tr_{H^s}P_0(Y_j^*Y_j)^{n_1}\right)^{1/n_1}.
\]
Hölder in $\xi$ and the Bloch trace formula
identify its $j$th integrated factor with
the zero-flux functional having the same
magnetic factor at every slice and no
electric factor. Let $\beta^{(j)}$ denote
that purely magnetic two-cochain. Monotone
convergence as $\delta\downarrow0$ gives
\[
 \mathcal Z(\beta)\le
                \prod_j\mathcal Z(\beta^{(j)})^{1/n_1}.
\]
Rotate the torus by an orientation-preserving
cyclic coordinate permutation. Zero flux
and the pairing are preserved, and
$\beta^{(j)}$ becomes purely electric for
the new time direction. The same inequality
then bounds $\mathcal Z(\beta^{(j)})$
by $\mathcal Z(0)$. Hence
$\mathcal Z(\beta)\le Z^{(0)}[1]$.

The zero-sector identity
\cite[(34)]{ReferenceResponse} identifies its
normalized curvature law with $\bar\mu$.
Since this law is supported on $F_N$,
the pairing with $g$ equals that with
$\operatorname{Pr}_{F_N}g$. Apply the
preceding inequality with the negative of
that projection and divide by the positive
zero-sector partition function. This proves
\eqref{eq:interacting-zero-flux-Gaussian}.
Exponential Markov applied to the two
signs and optimized at parameter $a$
gives the stated tail bound.
\end{proof}


\section{The terminal measure}
\label{sec:terminal-partition}

We first put the curvature Gaussian in coordinates
adapted to all block steps. The energy splits into
the terminal photon energy and independent unit
Gaussian fluctuation coordinates. Iterating the exact
density identity in these coordinates then gives
the terminal fermion integral on the small-field set.

\subsection{The downward photon coordinates}
\label{subsec:term-photon-tower}

Fix a flow side and cutoff as in
\cite[Proposition~6.1]{CoverTuning}. Put $k=k_*(R)$ and
$n=R/\eta=L^k\ell_*$.
On the unit torus $T_j^0=\Z^3/\ell_j\Z^3$ of side
$\ell_j=n/L^j$ use the
fluctuation space $V_j$ and its fixed orthonormal
basis, and the maps $A_j^{\min},\mathrm{Sq}_j$ of
the photon construction \cite[\S3.2--3.3]{MaxwellBlock}.
The map $A_j^{\min}$ is the constrained photon minimizer,
and $\mathrm{Sq}_j$ is the fluctuation map formed from the
positive covariance square root and the coordinate embedding
into constrained bond fields. We write $\mu_I$ for the standard Gaussian
on $V_j$ and
$\mathrm{Box}_j=\{\zeta_j:|\zeta_j(\Upsilon)|\le p_{0,j}
\text{ for every basis coordinate }\Upsilon\}$.
For a real terminal unit bond field $A$ and
$\zeta=(\zeta_j)_{j<k}$ define
\begin{equation}\label{eq:term-downward-recursion}
 A_k^\downarrow=A,\qquad
 A_j^\downarrow=A_j^{\min}(A_{j+1}^\downarrow)
                         +\mathrm{Sq}_j\zeta_j,\qquad j=k-1,\ldots,0 .
\end{equation}
The last field is denoted $A_0^\downarrow(A,\zeta)$.
These fields are linear in their data.

For any fine unit field $A_0$ define its block fields by
$A_j=(\mathbf Q^{(j)}A_0)^{[j]}$, where $\mathbf Q^{(j)}$
is the $j$-fold photon average on the original unit torus and
$B^{[j]}(b)=L^{j/2}B(L^jb)$.
The block-averaging/scaling identities give
$(A_{j+1})_L=\mathbf QA_j$.
Let $\mathrm{Ax}_k$ be the fields $A_0$ for which
$\tau A_j=0$ for $j<k$.
For each unit torus put
$V_{\mathrm t}=\delta C^2$, the real coclosed bond
fields with mean zero.
On the fine torus let $F_N=dC^1(n)$ with its
counting inner product and standard Gaussian
$\gamma_N$. The slice is
\[
 \mathrm{Sl}_k=\{A_0\in\mathrm{Ax}_k:A_k\in V_{\mathrm t}(T_k^0)\}.
\]

We will use the elementary cochain identities
\begin{equation}\label{eq:term-unit-Hodge}
 C^1=dC^0\oplus\{\text{constant bond fields}\}\oplus\delta C^2,
       \qquad \ker d=dC^0\oplus\{\text{constant bond fields}\}.
\end{equation}
The sum is orthogonal. Indeed $dd=0$, so $dC^0$
and $\delta C^2$ are orthogonal, and their common
orthogonal complement is $\ker(d\delta+\delta d)$.
Commuting lattice differences make this operator
the scalar Laplacian in each component.
Its kernel consists of constants: the energy identity
forces every nearest-neighbor difference to vanish.
This proves \eqref{eq:term-unit-Hodge}.
Consequently $d$ is injective on $V_{\mathrm t}$ and,
on a unit torus with $\ell^3$ sites,
$\dim V_{\mathrm t}=\dim dC^1=2\ell^3-2$.

\begin{proposition}\label{prop:term-photon-tower}
The map $(A,\zeta)\mapsto A_0^\downarrow(A,\zeta)$
is a linear bijection from
$C^1(T_k^0)\times\prod_{j<k}V_j$ to $\mathrm{Ax}_k$.
The block fields of its image are exactly
$A_j^\downarrow$.
The coboundary is a linear bijection
$d:\mathrm{Sl}_k\to F_N$.
Thus curvature has coordinates
\begin{equation}\label{eq:term-curvature-coordinates}
 \mathcal J_{\mathrm{ph}}:
 F\longmapsto(A_*;\zeta_0,\ldots,\zeta_{k-1}),\qquad
 F=dA_0^\downarrow(A_*,\zeta),
\end{equation}
where $A_*\in V_{\mathrm t}(T_k^0)$.
They satisfy the exact energy identity
\begin{equation}\label{eq:term-tower-energy}
 \|F\|^2=\|dH_k^{\mathrm{ph}}A_*\|_2^2
                                  +\sum_{j<k}\|\zeta_j\|^2 .
\end{equation}
The image of $\gamma_N$ under
$\mathcal J_{\mathrm{ph}}$ is
\begin{equation}\label{eq:term-tower-Gaussian}
                  \dd\mathsf m_*(A_*)\prod_{j<k}\dd\mu_I(\zeta_j),
\end{equation}
where the terminal centered Gaussian on
$V_{\mathrm t}(T_k^0)$ is nondegenerate and has density
proportional to
\[
                  \exp(-\|dH_k^{\mathrm{ph}}A\|_2^2/2).
\]

There are orthonormal vectors
$\phi_{j,\Upsilon}\in F_N$ such that
$\zeta_j(\Upsilon)=\langle\phi_{j,\Upsilon},F\rangle$.
For each terminal unit plaquette $p$ there is
$\phi_p\in F_N$ with
\begin{equation}\label{eq:term-tower-row-norms}
 (dA_*)(p)=\langle\phi_p,F\rangle,\qquad \|\phi_p\|\le1.
\end{equation}
There are $3\ell_*^3$ such plaquettes, and
\begin{equation}\label{eq:term-tower-coordinate-count}
 \dim V_j=(2-2L^{-3})\ell_j^3
                       \le2\ell_*^3L^{3(k-j)} .
\end{equation}
\end{proposition}

\begin{proof}
The one-step minimizer and fluctuation maps obey
$\mathbf QA_j^{\min}(A')=A'_L$, $\tau A_j^{\min}(A')=0$,
$\mathbf Q\mathrm{Sq}_j=0$ and $\tau\mathrm{Sq}_j=0$.
Apply $\mathbf Q$ to \eqref{eq:term-downward-recursion};
the injective scaling map then identifies the next
block field with $A_{j+1}^\downarrow$.
Starting from $A_0=A_0^\downarrow$ proves this at all
levels, including $\tau A_j=0$.
Conversely, for $A_0\in\mathrm{Ax}_k$ its difference
$A_j-A_j^{\min}(A_{j+1})$ is annihilated by $\mathbf Q$ and
$\tau$. The coordinate bijection
\cite[(21)]{MaxwellBlock} and the invertible
square root give exactly one $\zeta_j$ with this
difference equal to $\mathrm{Sq}_j\zeta_j$.
Downward recursion recovers $A_0$.
This proves the first bijection and its inverse.

The orthogonality
\cite[(28)]{MaxwellBlock}
and the Gaussian fluctuation decomposition give
at each level
\[
 \|dH_j^{\mathrm{ph}}A_j\|_2^2
       =\|dH_{j+1}^{\mathrm{ph}}A_{j+1}\|_2^2+\|\zeta_j\|^2 .
\]
Summing, and using $H_0^{\mathrm{ph}}=I$, proves
\eqref{eq:term-tower-energy}.
The dimension of the slice is
\[
 2\ell_k^3-2+\sum_{j<k}(2\ell_j^3-2\ell_{j+1}^3)
                                                    =2n^3-2=\dim F_N.
\]
If $dA_0=0$ on the slice, the energy identity forces
every $\zeta_j=0$ and $dH_k^{\mathrm{ph}}A_k=0$.
Since $\mathbf Q_kH_k^{\mathrm{ph}}=I$ and
$d\mathbf Q_k=\mathbf Q_k^{(2)}d$,
this implies $dA_k=0$.
The injectivity of $d$ on $V_{\mathrm t}$ gives
$A_k=0$, whence $A_0=0$ by the downward bijection.
Equality of dimensions therefore proves that
$d:\mathrm{Sl}_k\to F_N$ is bijective and gives
\eqref{eq:term-curvature-coordinates}.

Equip the coordinate space with the scalar product
\[
 \langle(A,\zeta),(A',\zeta')\rangle_{\mathrm{ph}}
       =\langle A,\Delta_kA'\rangle
                                      +\sum_{j<k}\langle\zeta_j,\zeta_j'\rangle.
\]
Its terminal part is positive definite:
$\langle A,\Delta_kA\rangle
=\|dH_k^{\mathrm{ph}}A\|_2^2\ge\|dA\|_2^2>0$
for nonzero $A\in V_{\mathrm t}$.
Here the block curvature contraction follows from
\cite[(5)]{MaxwellBlock}.
The energy identity and polarization make
$\mathcal J_{\mathrm{ph}}$ a linear isometry for
this scalar product.
It maps orthonormal bases to orthonormal bases;
the Jacobian is therefore $1$ for their Euclidean
volume measures. These are the measures in the
normalized standard Gaussians.
In an orthogonal direct sum, that volume measure
and the Gaussian density factor into those of
the summands. This proves
\eqref{eq:term-tower-Gaussian}, including its normalization.
The volume measure for the terminal scalar product
differs from the ordinary Lebesgue measure on
$V_{\mathrm t}$ by one constant factor.
Thus its Gaussian has exactly the stated density.
Its quadratic form is positive definite, so it is
nondegenerate.

For a basis vector $\Upsilon$ in $V_j$, let
$\mathbf e_{j,\Upsilon}$ be the coordinate vector
with that single nonzero entry.
Its inverse image under the isometry is
$\phi_{j,\Upsilon}$.
The coordinate vectors are orthonormal, proving
orthonormality and the required coordinate formula.
The terminal curvature functional satisfies
\[
 |(dA_*)(p)|
 \le\|dA_*\|_2
   =\|\mathbf Q_k^{(2)}dH_k^{\mathrm{ph}}A_*\|_2
 \le\|dH_k^{\mathrm{ph}}A_*\|_2\le\|F\|.
\]
It is consequently represented by a vector of norm
at most $1$ in the finite Euclidean space $F_N$,
giving \eqref{eq:term-tower-row-norms}.
Finally \cite[(19)]{MaxwellBlock} and
$\ell_j=L^{k-j}\ell_*$ give
\eqref{eq:term-tower-coordinate-count}.
\end{proof}


If a terminal field $A^c$ is closed, linearity gives
the difference of downward fields as
$A_0^\downarrow(A^c,0)$.
For a closed input the background identity
\cite[(30)]{MaxwellBlock} gives zero
curvature energy at the preceding level. Applying
the average shows that its block minimizer is closed.
The identity $\mathbf QA_j^{\min}(A')=A'_L$ and
preservation of means by the block average give
$\operatorname{mean}(A_j^{\min}(A'))=
L^{-1/2}\operatorname{mean}(A')$.
Its final mean is $L^{-k/2}\operatorname{mean}(A^c)$.
Equation~\eqref{eq:term-unit-Hodge} therefore gives
\begin{equation}\label{eq:term-closed-terminal-shift}
 A_0^\downarrow(A+A^c,\zeta)
    =A_0^\downarrow(A,\zeta)
        +L^{-k/2}\operatorname{mean}(A^c)+d\lambda .
\end{equation}
In particular adding a coarse gradient adds a fine
gradient. For the harmonic twist
$c_{\theta,\mu}=\theta_\mu/(e_*\ell_*)$,
$L^{-k/2}c_\theta=H_\theta$, where
$H_{\theta,\mu}=\theta_\mu/(qn)$.
Thus its downward shift is $H_\theta+d\lambda$.

\subsection{The terminal density with an external connection}
\label{subsec:term-density}

Let $(m_{\mathrm{in}},\varepsilon_{\mathrm{in}}^{\mathrm v})$
be a controlled real pair on $R$, with the trajectory
of \cite[\S7.1]{CoverTuning}. Thus its initial activity is zero,
the actual block recursion exists through $k=k_*(R)$, and
$|m_j|\le w_j^{\rm m}=e_j^{3/4-8\varepsilon}$ for $j\le k$.
Its activity satisfies
$\|E_j\|_j\le\|\mathscr E_j\|_j^x\le
4e_j^{1/4-6\varepsilon}\le w_j^{\rm E}
=e_j^{1/4-7\varepsilon}$
by \cite[Proposition~7.1]{CoverTuning}.
We retain the block field domains and the weighted polymer
norms of \cite[\S2.2--2.3]{NormalizedContraction} and \cite[\S2.3, \S5.1]{NonlinearStep}. The fine mesh at level
$j$ is $\eta_j=L^{-j}$. In particular,
\begin{align*}
 \mathrm{Dom}_j
  &=\{A\text{ real unit bond field}:
           H_j^{\rm ph}A\in\tfrac14\widetilde R_j\},\\*
 \mathfrak X_j
  &=\{\mathscr C:\|\mathscr C\|_\infty<e_j^{-1+2\varepsilon},\
                 \|\partial\mathscr C\|_\infty<e_j^{-1+3\varepsilon}\}.
\end{align*}
The domain $\widetilde R_j$ allows a global complex small field
and local real-gauge representatives with the three bounds
$e_j^{-3/4+i\varepsilon}$, $i=1,2,3$, in the supremum,
first-difference and H\"older first-difference norms;
its precise enlarged cubes are those of
\cite[\S6]{MaxwellBlock}.
An external connection is a complex fine unit bond
field $a$ satisfying
$a^{[j]}\in\mathfrak X_j$ for every $j\le k$;
here $a^{[j]}(b)=L^{j/2}a(L^jb)$.
Write $\mathscr C_*=a^{[k]}$.
For $A\in\mathrm{Dom}_k$ put
$\mathscr A=H_k^{\mathrm{ph}}A$,
$B=\mathscr A+\mathscr C_*$, and insert the fermion
fields $H_k^\#(B)\Psi$ and their barred partners.
The notation $H_k^\#$ includes $H_k\Psi$ and its covariant
H\"older differences, together with
$\bar H_k\bar\Psi$ and their barred differences;
the operators are defined in
\eqref{eq:interacting-terminal-operators}.
The kernel $\mathscr E_k(X;\mathscr A,\mathscr C,\psi,\chi)$
is the resulting Grassmann activity on an $M$-polymer $X$,
and $\mathfrak P_k(X,\mathscr C)$ is its external scalar family.
They are the actual recursively constructed families in
\cite[(7)]{NonlinearStep} and \cite[(55)--(57)]{CoverTuning}.
The terminal function is
\begin{align}
 Y_*^x(A;\mathscr C_*)
   ={}&\det D_k(0)^{-1}
        \exp\left(-\varepsilon_k^{\mathrm v}\ell_*^3
                       +\sum_X\mathfrak P_k(X,\mathscr C_*)\right)
       \nonumber\\
   &\times\int\mathrm D\Psi\,
      \exp\left(-\langle\bar\Psi,D_k(B)\Psi\rangle
                -m_k\langle\bar\psi,\psi\rangle
                       +\sum_X\mathscr E_k(X;\mathscr A,\mathscr C_*,
                                                        \psi,\chi)\right).
                                      \label{eq:term-Y-density}
\end{align}
Since $k\ge1$, this Berezin integral has $4\ell_*^3$
pairs with weight $1$. The inserted chamber fields
and kernels retain their fine weight $\eta_k^3/48$.
The inverse and field maps of \cite[Theorem~A and \S4.2]{FermionBlock},
together with the activity domains and exact density of
\cite[\S2.3 and \S3]{NonlinearStep}, make every term well defined.
At zero connection, the positive-reference-mass inverse formula in
\cite[Proposition~3.1]{NonlinearStep}
gives $\det D_k(0)\ne0$.
At zero external slot,
$\mathfrak P_k(X,0)=0$ and $\mathscr E_k|_0=E_k$;
we write $Y_*(A)=Y_*^x(A;0)$.

For the original fine unit field let
\begin{equation}\label{eq:term-fine-weight}
 w_N(B)=
 \exp(-\varepsilon_{\mathrm{in}}^{\mathrm v}n^3)
       \frac{\det(D^{\mathrm{ch}}[qB]+\bar m\eta+m_{\mathrm{in}})}
                              {\det(D^{\mathrm{ch}}[0]+\bar m\eta)} .
\end{equation}
The matrices have size $192n^3$.
The bracket argument of $D^{\mathrm{ch}}$ is the bond phase,
as in \cite[\S1.1, (6)]{PositiveProbe}; an unabsorbed unit
connection $B$ therefore has argument $qB$.
The denominator is nonzero by the free positive-mass
coercivity estimate; no invertibility of the numerator
is required.
The weighted Gaussian Berezin identity gives
the initial integrated density exactly as $w_N(A_0+a)$.
In particular all initial measure and vacuum factors
are included in \eqref{eq:term-fine-weight}.

This fine weight is entire in $B$ and is invariant
under real gauge shifts.
Its entries involve only the four chamber neighbors,
with link phases in the fourth neighbor.
Multiplication by $\exp(-\mathrm i q\lambda(c))$ on
all chamber components over $c$ conjugates the operator
at $B$ to the operator at $B+d\lambda$: the endpoint
phases cancel the gradient in the fourth hopping phase,
and the other hops stay over the same $c$.
The determinants therefore agree.
The same stencil bounds every row sum at $B_0+a$,
with $B_0$ real, by
$|\boldsymbol\mu|+41+6\exp(q\|\operatorname{Im}a\|_\infty)$,
where $\boldsymbol\mu=\bar m\eta+m_{\mathrm{in}}$.
Indeed the three internal hopping block norms sum
with the diagonal to less than $|\boldsymbol\mu|+41$,
and the fourth has row norm at most
$6\exp(q\|\operatorname{Im}a\|_\infty)$.
The determinant expansion bounds its modulus by
the product of the absolute row sums: its permutation
sum is a subsum of the sum over all column choices.
Consequently
\begin{equation}\label{eq:term-finite-density-bound}
 |w_N(B_0+a)|\le B_N(a):=
 \frac{\exp(-\varepsilon_{\mathrm{in}}^{\mathrm v}n^3)}
            {|\det(D^{\mathrm{ch}}[0]+\bar m\eta)|}
       \big(|\boldsymbol\mu|+41
                     +6\exp(q\|\operatorname{Im}a\|_\infty)\big)^{192n^3}
                                                        <\infty .
\end{equation}
This bound is used for integrability at a fixed cutoff.

Let $A(F)\in V_{\mathrm t}$ be the transverse fine
potential of $F\in F_N$, and define the finite complex
measure
\[
 \dd\nu^{\theta,a}(F)=
               w_N(A(F)+H_\theta+a)\dd\gamma_N(F).
\]
The small-field event and the terminal indicator are
\begin{equation}\label{eq:term-small-field-event}
 \mathfrak G=\{F:|\zeta_j(\Upsilon)|\le p_{0,j}\ (j<k),
                                 \ \|dA_*\|_\infty\le p_*\},\qquad
 \chi_*(A)=\mathbf1[\|dA\|_\infty\le p_*].
\end{equation}
Here $A_*,\zeta$ are the actual coordinates
\eqref{eq:term-curvature-coordinates}.
Products with $\chi_*$ are extended by zero off
its support.

\begin{theorem}\label{thm:term-density}
For every $A\in\mathrm{Dom}_k$,
\begin{equation}\label{eq:term-iterated-density}
 Y_*^x(A;\mathscr C_*)
      =\int\prod_{j<k}
                    \big(\mathbf1_{\mathrm{Box}_j}(\zeta_j)
                                             \dd\mu_I(\zeta_j)\big)
                       w_N(A_0^\downarrow(A,\zeta)+a).
\end{equation}
The integrand is continuous on the compact product
of boxes.
If $A,A+d\lambda\in\mathrm{Dom}_k$, with $\lambda$
real, their terminal functions agree.
For every bounded Borel function $\Phi$ on
$V_{\mathrm t}(T_k^0)$ and every real twist $\theta$,
\begin{equation}\label{eq:term-Borel-observable}
 \nu^{\theta,a}(\mathbf1_{\mathfrak G}\Phi(A_*))
    =\int\dd\mathsf m_*(A)\,
                        \Phi(A)\chi_*(A)Y_*^x(A+c_\theta;\mathscr C_*).
\end{equation}
Its right-hand density is bounded and Borel, and
$Y_*^x(A+c_\theta;\mathscr C_*)$ is jointly continuous
in $(\theta,A)$ on $\{\chi_*=1\}$.
\end{theorem}

\begin{proof}
The photon domain conditions
\cite[(43)]{MaxwellBlock} hold at each
running scale by the fixed charge choice.
\cite[Theorem~B]{MaxwellBlock} gives
two consequences used here: a real unit field with
curvature bounded by $p_j$ lies in $\mathrm{Dom}_j$;
and a field in $\mathrm{Dom}_{j+1}$ descends to
$\mathrm{Dom}_j$ after addition of any box fluctuation.
Downward induction therefore puts every field of
\eqref{eq:term-downward-recursion} in its respective
domain whenever its top field lies in $\mathrm{Dom}_k$ and
all fluctuations lie in their boxes.

Let $y_j^x(A;a^{[j]})$ be the integrated density
at level $j$ along this trajectory.
The exact integrated step \cite[(9)]{NonlinearStep} gives
\[
 y_{j+1}^x(A';a^{[j+1]})
     =\int_{\mathrm{Box}_j}\dd\mu_I(\zeta_j)\,
                   y_j^x(A_j^{\min}(A')+\mathrm{Sq}_j\zeta_j;
                                                               a^{[j]}).
\]
Its external slot is the required one because
$(a^{[j+1]})_L=a^{[j]}$, directly from the two
scaling factors.
Its coupling maps are the actual trajectory maps,
and its vacuum includes the box normalization.
The first such identity integrates the full
$192n^3$ weighted chamber pairs into
$4\ell_1^3$ coarse pairs.

For a top field in $\mathrm{Dom}_j$, the function
$w_N(A_0^\downarrow(A,\zeta)+a)$ is continuous in
the finitely many fluctuation coordinates.
Every finite product of boxes is compact,
and \eqref{eq:term-finite-density-bound} bounds it.
Thus all repeated integrals are absolutely integrable.
Starting with the initial identity
$y_0^x(A_0;a)=w_N(A_0+a)$, use the one-step formula,
the downward recursion and Fubini at each level.
This proves \eqref{eq:term-iterated-density} at $j=k$,
with its unnormalized factors
$\mathbf1_{\mathrm{Box}_j}\dd\mu_I$.
The continuity assertion follows from the same
finite-dimensional composition.
Equation~\eqref{eq:term-closed-terminal-shift}
turns a coarse real gradient into a fine gradient;
gauge invariance of $w_N$ inside this integral
then proves gauge invariance of $Y_*^x$.

To integrate against an arbitrary bounded Borel
observable, first express the fine density in the
curvature coordinates.
The slice potential $A_0^\downarrow(A_*,\zeta)$
has mean zero: its block average has the mean of
the fine field, and its top field has mean zero.
Its difference from $A(F)$ is closed and has mean
zero, hence is a gradient by
\eqref{eq:term-unit-Hodge}.
The harmonic shift identity gives
\[
 w_N(A(F)+H_\theta+a)
             =w_N(A_0^\downarrow(A_*+c_\theta,\zeta)+a).
\]
Under the tower isometry the Gaussian becomes
\eqref{eq:term-tower-Gaussian}, and the event
$\mathfrak G$ has indicator
\[
              \chi_*(A_*)\prod_{j<k}\mathbf1_{\mathrm{Box}_j}(\zeta_j).
\]
The integrand is a bounded Borel function with bound
$\|\Phi\|_\infty B_N(a)$.
The change-of-variables identity for its image
measure and Fubini therefore give
\begin{align*}
 \nu^{\theta,a}(\mathbf1_{\mathfrak G}\Phi(A_*))
   =\int\dd\mathsf m_*(A)\,\Phi(A)\chi_*(A)
      \int\prod_{j<k}
                      \big(\mathbf1_{\mathrm{Box}_j}\dd\mu_I(\zeta_j)\big)
                         w_N(A_0^\downarrow(A+c_\theta,\zeta)+a).
\end{align*}
Where $\chi_*(A)=1$, the field $A+c_\theta$ has
the same curvature bound and belongs to
$\mathrm{Dom}_k$. The inner integral is
$Y_*^x(A+c_\theta;\mathscr C_*)$ by
\eqref{eq:term-iterated-density}, proving
\eqref{eq:term-Borel-observable}.

The inner integral defines a function on all real
$(\theta,A)$, whether or not $A+c_\theta\in\mathrm{Dom}_k$.
Its modulus is at most $B_N(a)$.
Its integrand is continuous and uniformly continuous
on a compact neighborhood of any $(\theta,A)$
times the compact box product. Integrating this
uniform bound proves joint continuity.
Multiplying by $\chi_*$ gives a bounded Borel
function, since $\{\chi_*=1\}$ is closed.
On that set it is the asserted terminal density.
\end{proof}

\subsection{The large-field complement}
\label{subsec:term-complement}

At zero external connection write $\nu^\theta=\nu^{\theta,0}$,
and let $\bar\nu$ be its average over $[-\pi,\pi]^3$
with measure $\dd\theta/(2\pi)^3$.
The real chamber operator has two adjoint spin blocks;
thus its determinant at every real lattice mass is
the squared modulus of the determinant of one block.
Consequently $\nu^\theta$ is positive.
The zero-sector identity
\cite[(40)]{ReferenceResponse} gives
$0<\bar\nu(1)<\infty$ and identifies
$\bar\mu=\bar\nu/\bar\nu(1)$ with the zero-flux
curvature probability law.

Use the fixed constant $C_0\ge1$ of
\cite[(25)]{QED3Integration}, recalled in
\cref{subsec:term-remainder}. It also bounds the terminal
field maps and Dirac form in
\cref{subsec:interacting-terminal-norms}. Set
\begin{equation}\label{eq:term-complement-errors}
 \varepsilon_3
 =6\ell_*^3\left(e^{-p_{0,*}/2}
                    +e^{-p_*^2/(2C_0^2)}\right),\qquad
 \varepsilon_4=6\ell_*^3e^{-p_*^2/(2C_0^2)}.
\end{equation}

\begin{proposition}\label{prop:term-complement}
Suppose $v_*\ge8$, the integer $p_0\ge1$ satisfies
$p_0v_*^{p_0-1}\ge28$, and $L\ge5$. Then, at every
real initial pair,
\begin{equation}\label{eq:term-complement-bound}
 \bar\mu(\mathfrak G^c)\le\varepsilon_3,\qquad
             \mathsf m_*(\chi_*=0)\le\varepsilon_4.
\end{equation}
For fixed $(\bar m,e,R)$, these bounds are independent
of the ultraviolet cutoff and of the initial pair.
\end{proposition}

\begin{proof}
The zero-flux domination estimate in
\cref{prop:interacting-zero-flux-domination}
and the rows of \eqref{eq:term-tower-row-norms} give
\[
 \bar\mu\{|\zeta_j(\Upsilon)|>p_{0,j}\}
                   \le2e^{-p_{0,j}^2/2},\qquad
 \bar\mu\{|(dA_*)(p)|>p_*\}\le2e^{-p_*^2/2}.
\]
Using the coordinate counts in
\eqref{eq:term-tower-coordinate-count}, the union bound is
\begin{equation}\label{eq:term-complement-union}
 \bar\mu(\mathfrak G^c)
 \le4\ell_*^3\sum_{i=1}^{k}L^{3i}
                  e^{-(v_*+i\log L/2)^{2p_0}/2}
                +6\ell_*^3e^{-p_*^2/2}.
\end{equation}
Here $i=k-j$ and $v_j=v_*+i\log L/2$.
The first two terms of the binomial expansion give
\[
 p_{0,j}=(v_*+i\log L/2)^{p_0}
       \ge p_{0,*}+28i\log L/2.
\]
Since $p_{0,j}\ge1$, its square is at least $p_{0,j}$.
It follows that
\[
 L^{3i}e^{-p_{0,j}^2/2}
                  \le e^{-p_{0,*}/2}L^{-4i}.
\]
Summing the geometric series, whose sum is
$(L^4-1)^{-1}\le1$, bounds
\eqref{eq:term-complement-union} by
$4\ell_*^3e^{-p_{0,*}/2}+6\ell_*^3e^{-p_*^2/2}$.
This is at most $\varepsilon_3$, since $C_0\ge1$.

The image of $\gamma_N$ under $F\mapsto A_*$ is
$\mathsf m_*$ by \eqref{eq:term-tower-Gaussian}.
For each plaquette the variable
$\langle\phi_p,F\rangle$ is Gaussian with variance
$\|\phi_p\|^2\le1$. Its two-sided exponential Markov
bound is $2e^{-p_*^2/2}$.
Summing over the $3\ell_*^3$ plaquettes proves
$\mathsf m_*(\chi_*=0)\le6\ell_*^3e^{-p_*^2/2}
\le\varepsilon_4$.
Finally, $\ell_*$ and $e_*$ are independent of $N$
after the mesh threshold, by
\cite[Proposition~6.1]{CoverTuning}; hence the displayed
errors have the asserted dependence.
\end{proof}

\section{Free pressure and harmonic twists}
\label{subsec:term-free-determinants}

The partition comparison requires an error in the free pressure
which depends on the physical size, rather than on the number of
lattice sites. It also requires a twist estimate for the operator
with a block penalty. The latter penalty removes the small-mass
singularity before the fine mesh is taken to zero.

For $0<m\le1$ let
$\mathfrak d_m(p)=\det\mathcal D_m(p)$, where
$\mathcal D_m$ is the $192$-dimensional chamber symbol in
\cite[\S2.2]{FermionBlock}, and set
\begin{equation}\label{eq:term-free-pressure}
 p_{\rm ch}^0(m)=
       \int_{[-\pi,\pi]^3}\log\mathfrak d_m(p)
                                      \frac{\dd p}{(2\pi)^3}.
\end{equation}
The spin-block identity and positive real part imply
$\mathfrak d_m(p)>0$ for real $p$.
In this subsection $|z|=|z|_\infty$ for $z\in\Z^3$.
Write $\kappa_{(0)}=\kappa_{\mathrm{strip}}$ from
\cite[Proposition~2.2]{FermionBlock}.

\begin{proposition}\label{prop:term-free-pressure}
There is a finite $C_6$, depending only on $\rho$, such that,
with
\begin{equation}\label{eq:term-epsilon-five}
 \varepsilon_5(x)=
 C_6\log\frac1{1-e^{-\kappa_{(0)}x}}\qquad(x>0),
\end{equation}
the free twisted determinant on a unit torus of side $n\ge1$
satisfies
\begin{equation}\label{eq:term-free-pressure-error}
 \left|\log\det(D^{\rm ch}[qH_\theta]+m)
                     -n^3p_{\rm ch}^0(m)\right|
                                 \le\varepsilon_5(mn).
\end{equation}
Here $H_{\theta,\mu}=\theta_\mu/(qn)$ and $\theta\in\R^3$.
For $m=\bar m\eta$ and $n=R/\eta$ define
\[
 z_{r,n}^{\rm free}
       =\int_{[-\pi,\pi]^3}
            \det(D^{\rm ch}[qH_\theta]+\bar m\eta)
                                         \frac{\dd\theta}{(2\pi)^3}.
\]
Then $z_{r,n}^{\rm free}>0$ and
\begin{equation}\label{eq:term-free-average-error}
 \left|\log z_{r,n}^{\rm free}
                 -n^3p_{\rm ch}^0(\bar m\eta)\right|
                              \le\varepsilon_5(\bar mR).
\end{equation}
\end{proposition}

\begin{proof}
Put $P_{n,\theta}=\{(2\pi j+\pi\mathbf1+\theta)/n:
j\in\{0,\ldots,n-1\}^3\}$. The antiperiodic plane waves
have momenta $p\in P_{n,0}$ and diagonalize the translation
part of the chamber operator. The external phase in the
fourth hop evaluates its symbol at $p+\theta/n\in P_{n,\theta}$,
whereas the internal hops are unchanged. Thus
\begin{equation}\label{eq:term-free-product}
 \det(D^{\rm ch}[qH_\theta]+m)
                   =\prod_{p\in P_{n,\theta}}\mathfrak d_m(p)>0.
\end{equation}
The phases contain $qH_\theta=\theta/n$, so the determinant
is independent of $q$.

Factor the determinant by the light--heavy decomposition:
\[
 \log\mathfrak d_m(p)=f_H(p)+f_S(p),\qquad
 f_H(p)=\log\det D_{\rm hh}(p),\quad
 f_S(p)=\log\det B_m(p).
\]
Both logarithms are real and smooth on real momenta.
Their derivatives, which suffice for contour shifting, are
\[
 \partial_\mu f_H=\tr(D_{\rm hh}^{-1}\partial_\mu D_{\rm hh}),
 \qquad
 \partial_\mu f_S=\tr(B_m^{-1}\partial_\mu B_m).
\]
The heavy block has dimension $188$. On the strip of
width $\theta_0$ its inverse has norm at most $2/g_{\rm h}$,
and the fourth-hop derivative has norm at most
$3(1+\rho)e^{\theta_0}=3(1+\rho)+g_{\rm h}/2$.
Hence
\[
 |\partial_\mu f_H|\le
 C_{\rm lH}:=188\left(1+\frac{6(1+\rho)}{g_{\rm h}}\right).
\]
For $z\ne0$ choose $|z_\mu|=|z|$. One integration by parts
and a shift of that contour by
$-\mathrm i\theta_0\operatorname{sgn}(z_\mu)$ give
\begin{equation}\label{eq:term-heavy-log-Fourier}
 |\widehat f_H(z)|\le C_{\rm lH}|z|^{-1}e^{-\theta_0|z|}.
\end{equation}
The derivative trace is holomorphic in a neighborhood of
the closed strip. Its periodicity cancels the vertical sides
of the contour rectangle; no choice of a complex logarithm
is needed.

For the light factor put $r_m(p')=(|p'|_2^2+m^2)^{1/2}$.
The inverse derivative bounds
\cite[(15)]{FermionBlock} hold on
$|\operatorname{Im}p_\nu|\le\kappa_{(0)}m$.
The bound $K_s$ for $B_m$ on the strip of width $\theta_0$
and Cauchy's formula on polydiscs of radius $\theta_0/2$
give
$\|\partial^\beta B_m\|
\le24K_s(2/\theta_0)^{|\beta|}$ for $1\le|\beta|\le4$.
Leibniz's rule in the derivative trace therefore gives
\begin{equation}\label{eq:term-light-log-derivatives}
 |\partial^\alpha\partial_\mu f_S(p'+\mathrm ip'')|
 \le C_{\rm lS}\,r_m(p')^{-1-|\alpha|}
 \le C_{\rm lS}|p'|_2^{-1-|\alpha|}
       \quad(|\alpha|\le3,\ p'\ne0),
\end{equation}
where one may take
\[
 C_{\rm lS}
 =\max_{|\alpha|\le3}96K_s
   \sum_{\beta\le\alpha}
    \binom{\alpha}{\beta}C_\beta(\rho)
    (2/\theta_0)^{|\alpha|-|\beta|+1}
    (3\pi^2+1)^{(|\alpha|-|\beta|)/2}.
\]
This constant is independent of $m$.
Shift the differentiated Fourier integral by
$-\mathrm i\kappa_{(0)}m\operatorname{sgn}(z_\mu)e_\mu$.
The periodic Fourier estimate proved in
\cite[Proposition~2.3]{FermionBlock}, with singular
order $1$ and three derivatives, bounds this integral by
$c_{\rm F}C_{\rm lS}(1+|z|)^{-2}$, where $c_{\rm F}$
is its constant at those orders. Thus
\begin{equation}\label{eq:term-light-log-Fourier}
 |\widehat f_S(z)|
    \le c_{\rm F}C_{\rm lS}|z|^{-1}(1+|z|)^{-2}
                          e^{-\kappa_{(0)}m|z|}.
\end{equation}
The strip is contained in the heavy strip since
$\kappa_{(0)}m\le\theta_0/2$.
The inverse is holomorphic near the closed strip, so
the contour shift is legitimate.

The Fourier coefficients of $f=f_H+f_S$ are summable.
For a summable periodic Fourier series, summing over
$P_{n,\theta}$ keeps precisely the indices $z=nw$:
\begin{equation}\label{eq:term-Fourier-sampling}
 \sum_{p\in P_{n,\theta}}f(p)
 =n^3\sum_{w\in\Z^3}\widehat f(nw)
                       (-1)^{w_1+w_2+w_3}e^{\mathrm i\theta\cdot w}.
\end{equation}
This follows by summing the geometric progression in each
coordinate; absolute convergence permits the interchange.
Its $w=0$ term is $n^3p_{\rm ch}^0(m)$.
There are $24t^2+2\le26t^2$ indices with $|w|=t$.
The contribution of \eqref{eq:term-light-log-Fourier} is at most
\[
 26c_{\rm F}C_{\rm lS}
                  \sum_{t\ge1}t^{-1}e^{-\kappa_{(0)}mnt}.
\]
For the heavy term use
$\kappa_{(0)}m\le\theta_0/2$ and
$x^2e^{-\theta_0x/2}\le16/(\exp(1)\theta_0)^2$.
Its contribution is at most
\[
 \frac{416C_{\rm lH}}{(\exp(1)\theta_0)^2}
                     \sum_{t\ge1}t^{-1}e^{-\kappa_{(0)}mnt}.
\]
Since $\sum_{t\ge1}e^{-xt}/t=\log(1-e^{-x})^{-1}$,
\eqref{eq:term-free-pressure-error} follows with
$C_6=26[c_{\rm F}C_{\rm lS}
          +16C_{\rm lH}/(\exp(1)\theta_0)^2]$.
Finally every twisted determinant lies between
$\exp(n^3p_{\rm ch}^0(m)\pm\varepsilon_5(mn))$.
Its harmonic average lies in the same positive interval.
Taking logarithms proves \eqref{eq:term-free-average-error}.
\end{proof}

We next retain the block penalty. Let $\mathfrak N\ge3$
be odd and put $\tau=\mathfrak N^{-1}$.
On the fine torus of side $\ell\in\N$ let
$\mathcal H_\theta$ be the chamber spinors with mesh $\tau$
and boundary multipliers $-e^{\mathrm i\theta_\mu}$,
where $\theta\in\R^3$.
The scalar product uses $\tau^3/48$.
The coarse space has four components per unit site and
counting scalar product. The average $Q$ takes the mean
over each unit block and its $48$ labels, and $Q^*$
extends a coarse field constantly on that block.
Thus $QQ^*=1$, $\|Q\|=\|Q^*\|=1$, and
$\Pi_{\rm bl}=Q^*Q$ is an orthogonal projection.
Fix $0<b_-\le b_+<\infty$; the block construction uses
$b_-=1-L^{-1}$ and $b_+=1$.
Let $D$ be the free chamber Dirac operator at mesh $\tau$
on $\mathcal H_\theta$. For $0<\mu\le1$ and $b_-\le b\le b_+$ set
\begin{equation}\label{eq:term-free-block-operators}
 T_\theta=D+\mu+b\Pi_{\rm bl},\qquad
 D_\theta^{\rm bl}=b-b^2QT_\theta^{-1}Q^*.
\end{equation}
All determinants below are determinants of linear maps.
The weighting of a scalar product changes their matrix
representations, but does not add a determinant factor.

\begin{proposition}\label{prop:term-block-twist}
The determinants of $T_\theta$, $D+\mu$, and
$D_\theta^{\rm bl}$ are positive, and
\begin{align}
 (D_\theta^{\rm bl})^{-1}
     &=b^{-1}+Q(D+\mu)^{-1}Q^*,
                                    \label{eq:term-free-block-inverse}\\
 \det D_\theta^{\rm bl}
     &=b^{4\ell^3}\frac{\det(D+\mu)}{\det T_\theta}.
                                    \label{eq:term-free-block-determinant}
\end{align}
There are $C_\Sigma<\infty$ and $0<\gamma_\Sigma\le1$,
depending only on $(b_-,b_+,\rho)$, such that
\begin{equation}\label{eq:term-penalty-twist-bound}
 \left|\log\frac{\det T_\theta}{\det T_0}\right|
                  \le C_\Sigma\ell^{-1}e^{-\gamma_\Sigma\ell}.
\end{equation}
For the level $k\ge1$ data
\[
             (\mathfrak N,\ell,\mu,b)=(L^k,\ell_*,\mu_*,b_k),
             \qquad c_{\theta,\mu}=\theta_\mu/(e_*\ell_*),
\]
put
$R_\theta^{\rm tw}=\det D_k(c_\theta)/\det D_k(0)$.
Then
\begin{equation}\label{eq:term-twist-determinant-identity}
 R_\theta^{\rm tw}
 =\frac{\det(D^{\rm ch}[qH_\theta]+\bar m\eta)}
              {\det(D^{\rm ch}[0]+\bar m\eta)}
                  \frac{\det T_0}{\det T_\theta}>0.
\end{equation}
\end{proposition}

\begin{proof}
The real part of $D+\mu$ is at least $\mu$.
The same holds for $T_\theta$, since its penalty is positive.
Both operators have two adjoint spin blocks; their determinants
are therefore positive. Put $G=(D+\mu)^{-1}$,
$S=QGQ^*$ and $Z=(b^{-1}+S)^{-1}$.
Their inverses exist since $\operatorname{Re}G\ge0$ and
$\operatorname{Re}(b^{-1}+S)\ge b^{-1}$.
Multiplication gives
$T_\theta^{-1}=G-GQ^*ZQG$ and
$D_\theta^{\rm bl}=Z$, proving
\eqref{eq:term-free-block-inverse}.
Sylvester's determinant identity, applied to
$-\!bQ$ and $T_\theta^{-1}Q^*$, gives
\eqref{eq:term-free-block-determinant}. Its right side is
positive.

We prove \eqref{eq:term-penalty-twist-bound} with constants
independent of $\tau$ and $\mu$.
Let
\[
 c_{\rm nz}=\min\left\{\frac\rho\pi,
             \frac{g_{\rm h}\rho^2}{49\sqrt3\,\pi^3}\right\}>0.
\]
For real $q'\in[-\pi,\pi]^3$,
\begin{equation}\label{eq:term-symbol-nozero}
             \|\mathcal D_0(q')u\|
                         \ge c_{\rm nz}|q'|_2\|u\|.
\end{equation}
Normalize $\|u\|=1$ and
write $u=Pu+\Pi_{\rm h}u$ for its light and heavy parts,
where $\Pi_{\rm h}=1-P$ is the cluster-heavy projection.
If $d=\|\mathcal D_0(q')u\|$, then
$g_{\rm h}\|\Pi_{\rm h}u\|^2\le d$.
The light compression is
$W(q')=\rho\sum_\mu(1-\cos q'_\mu)+
\mathrm i\gamma\cdot\sin q'$; its squared modulus is
at least $(2\rho/\pi)^2|q'|_2^2$.
Also $J^*\mathcal D_0(0)=0$, and the fourth-hop
difference has norm at most $6|q'|_2$.
Projecting $\mathcal D_0(q')u$ onto the light space gives
\[
 d\ge |q'|_2\left(\frac{2\rho}{\pi}-7\|\Pi_{\rm h}u\|\right).
\]
If $\|\Pi_{\rm h}u\|\le\rho/(7\pi)$ this is at least
$(\rho/\pi)|q'|_2$. Otherwise
$d\ge g_{\rm h}\rho^2/(49\pi^2)$; since
$|q'|_2\le\sqrt3\pi$, this proves
\eqref{eq:term-symbol-nozero}.

On one unit cell use the Bloch space
$\mathcal H(p)=\{f(c+e_\mu,g)=e^{\mathrm ip_\mu}f(c,g)\}$.
Its modes are $p_j=p+2\pi j$, with
$j\in I_{\mathfrak N}^3$,
$I_{\mathfrak N}=\{-(\mathfrak N-1)/2,\ldots,
(\mathfrak N-1)/2\}$.
The map
$\{u_j\}\mapsto\sum_j e^{\mathrm ip_j\cdot c}u_j(g)$
is unitary for the one-cell scalar product $\tau^3/48$.
It diagonalizes $D$ with blocks
$X_j=\tau^{-1}\mathcal D_0(\tau p_j)$.
Define
\[
 u_1(t)=\mathfrak N^{-1}\sum_{s\in I_{\mathfrak N}}
                                      e^{-\mathrm its/\mathfrak N},
 \qquad u(p)=\prod_\mu u_1(p_\mu).
\]
The block projection in this representation has kernel
$u(p_j)u(p_{j'})P$. This follows by averaging the plane
waves on the centered unit block; discrete Fourier
orthogonality expands the constant block extension.
For $p\in[-\pi,\pi]^3$ one has $\tau p_j\in[-\pi,\pi]^3$
and $|p_j|_2\ge\pi|j|$ when $j\ne0$.
Hence \eqref{eq:term-symbol-nozero} gives
$\|X_ju\|\ge c_{\rm nz}|p_j|_2\|u\|$.

The inverse on $\mathcal H(p)$ is uniformly bounded.
To see this, let $\|f\|=1$ and
$d=\|(D+\mu+b\Pi_{\rm bl})f\|\le1$.
The real-part estimate gives
\[
 \mu\le d,\qquad
 \|\Pi_{\rm bl}f\|\le(d/b_-)^{1/2},\qquad
 \|\Pi_{\rm h}f\|\le(d/(3g_{\rm h}))^{1/2},
 \qquad \|Df\|\le(2+\sqrt{b_+})\sqrt d.
\]
Separate the mode $j=0$ from $f^{\rm r}$, the remaining
modes. The bound on $X_j$ gives
$\|f^{\rm r}\|\le\|Df\|/(\pi c_{\rm nz})$.
Moreover
$u_1(t)=\sin(t/2)/[\mathfrak N\sin(t/(2\mathfrak N))]
\ge2/\pi$ for $|t|\le\pi$, with value $1$ at zero.
Thus the block projection of the zero mode controls
its light part by the factor $8/\pi^3$.
Using $\|\Pi_{\rm bl}\|=1$ and subtracting the other modes,
\[
 \sqrt{d/b_-}\ge
 \frac8{\pi^3}(1-\|\Pi_{\rm h}f\|)
                     -(1+8/\pi^3)\|f^{\rm r}\|.
\]
It follows that $8/\pi^3\le K_{\rm iv}\sqrt d$, where
\[
 K_{\rm iv}=b_-^{-1/2}+\frac8{\pi^3}(3g_{\rm h})^{-1/2}
       +(1+8/\pi^3)\frac{2+\sqrt{b_+}}{\pi c_{\rm nz}}.
\]
Consequently the inverse norm is at most
$C_{\rm iv}=\max\{1,(\pi^3K_{\rm iv}/8)^2\}$.

For complex momenta use the periodic gauge
$f(c,g)\mapsto e^{\mathrm ip\cdot c}f(c,g)$ to work
on the fixed periodic one-cell space. Let $D(p)$
and $\Pi(p)$ be the resulting operators, so
$\mathcal T(p)=D(p)+\mu+b\Pi(p)$.
The fourth hopping phase is
$e^{\mathrm i\tau p\cdot ge_1}$, and the block kernel is
$e^{-\mathrm ip\cdot(c-c')}$. A shift of $p_\nu$
by $2\pi$ conjugates these operators by multiplication
by $e^{2\pi\mathrm ic_\nu}$.
Thus their determinant and derivative traces are periodic.
For $p=p'+\mathrm it e_\nu$, $|t|\le1$, and $1\le j\le4$,
direct differentiation gives
\begin{equation}\label{eq:term-Bloch-derivatives}
 \|\partial_\nu^jD(p)\|\le9\tau^{j-1},\qquad
 \|\partial_\nu^j\Pi(p)\|\le\exp(1),\qquad
 \operatorname{rank}\partial_\nu^j\Pi(p)\le4(j+1).
\end{equation}
The last bound follows by expanding $(c_\nu-c'_\nu)^j$
into $j+1$ products of a function of $c$ and a function
of $c'$, each with four spin components.
The norm estimates use $|c_\nu-c'_\nu|<1$ in one block,
$\tau\le1/3$, and $\|M_3\|\le6$.
The difference from the real momentum has norm at most
$(6+b_+)|t|e^{|t|}$. With
\[
 \gamma_\Sigma=
 \min\{1,[2\exp(1)(6+b_+)C_{\rm iv}]^{-1}\},
\]
a Neumann series gives
$\|\mathcal T(p)^{-1}\|\le2C_{\rm iv}$ on this strip.
Its singular values are at most twice those at $p'$.

The growing dimension of the fine cell cannot be used
in a trace bound. Instead write $G(p)=\mathcal T(p)^{-1}$
and compare at real $p'$ with $(D(p')+1)^{-1}$.
The resolvent identity, the bound just proved, and the
singular-value multiplication inequality
\cite[slides~3--9]{PongeSchatten2025} imply
\[
 s_j(G(p))\le C_{\rm sv}s_j((D(p')+1)^{-1}),\qquad
 C_{\rm sv}=2[1+(1+b_+)C_{\rm iv}].
\]
The comparison operator is the direct sum of
$(X_j+1)^{-1}$. Since $\operatorname{Re}X_j\ge0$,
each of its $192$ singular values is at most
$(1+c_{\rm nz}^2|p'_j|_2^2)^{-1/2}$.
Set
\[
 S_4=\sum_{j\in\Z^3\setminus\{0\}}|j|^{-4}
        =4\pi^2+\pi^4/45,\quad
 C_{\rm uv}=(1+3\pi^2c_{\rm nz}^2)^{1/2},\quad
 S_{\rm jf}=1+S_4/(\pi c_{\rm nz})^4.
\]
For $1\le s\le4$ we obtain
\begin{equation}\label{eq:term-Bloch-Schatten}
 \tau^{4-s}\|G(p)\|_{\mathfrak S_s}^{\,s}
           \le192C_{\rm sv}^{\,s}C_{\rm uv}^{\,4-s}S_{\rm jf}.
\end{equation}
Indeed
$\tau^2(1+c_{\rm nz}^2|p'_j|_2^2)\le C_{\rm uv}^2$,
so multiplication by $\tau^{4-s}$ makes each summand
a multiple of $(1+c_{\rm nz}^2|p'_j|_2^2)^{-2}$.
The nonzero indices are bounded by
$(\pi c_{\rm nz}|j|)^{-4}$ and have the summable
majorant displayed above.

Let $F(p)=\log\det\mathcal T(p)$ on real momenta.
The fourth derivative is the trace expression
\begin{align}
 F_4(p)={}&\tr(GT_4)-4\tr(GT_1GT_3)
                  -3\tr(GT_2GT_2)\nonumber\\
          &+12\tr(GT_1GT_1GT_2)-6\tr((GT_1)^4),
                                      \label{eq:term-Bloch-fourth}
\end{align}
where $T_j=\partial_\nu^j\mathcal T$.
This follows by differentiating
$\partial_\nu F=\tr(GT_1)$ and
$\partial_\nu G=-GT_1G$, using trace cyclicity.
The expression is holomorphic near the closed strip
in the $\nu$ coordinate.
Each trace in \eqref{eq:term-Bloch-fourth} has $s\le4$
inverse factors and derivatives of total order four.
Expand every $T_j$ as $D_j+b\Pi_j$.
For the term with only $D_j$, Schatten Hölder and
\eqref{eq:term-Bloch-derivatives} give the bound
$9^s\tau^{4-s}\|G\|_{\mathfrak S_s}^{\,s}$.
Equation~\eqref{eq:term-Bloch-Schatten} bounds this by
$9^4\,192C_{\rm sv}^4C_{\rm uv}^3S_{\rm jf}$.
In each other term estimate one $b\Pi_j$ in trace norm:
its rank bound gives $\|b\Pi_j\|_{\mathfrak S_1}
\le20\exp(1)b_+$. The other factors have operator
norms at most $2C_{\rm iv}$ and
$\max\{9,\exp(1)b_+\}$.
There are at most $15$ such terms in any trace.
Since the five numerical coefficients have absolute
sum $26$, this proves $|F_4(p)|\le C_{\rm Flg}$, where
\[
 C_{\rm Flg}=26\left[
       9^4\,192C_{\rm sv}^4C_{\rm uv}^3S_{\rm jf}
       +300\exp(1)b_+(2C_{\rm iv})^4
                           \max\{9,\exp(1)b_+\}^3\right].
\]
All Schatten applications are on a finite-dimensional
space, with $s$ factors of exponent $s$ and reciprocal
exponents summing to one
\cite[slides~54--56]{PongeSchatten2025}.

For $z\ne0$ choose $|z_\nu|=|z|$.
Four integrations by parts and a contour shift by
$-\mathrm i\gamma_\Sigma\operatorname{sgn}(z_\nu)e_\nu$
now give
$|\widehat F(z)|\le C_{\rm Flg}|z|^{-4}
e^{-\gamma_\Sigma|z|}$.
The trace expression is periodic by the preceding
conjugacy and is holomorphic near the closed strip.
Decomposition into unit-cell Bloch spaces gives
$\log\det T_\theta=\sum_{p\in P_{\ell,\theta}}F(p)$.
Using \eqref{eq:term-Fourier-sampling} at $\theta$ and
at zero, and subtracting, proves
\[
 \left|\log\frac{\det T_\theta}{\det T_0}\right|
 \le2\ell^3\sum_{w\ne0}C_{\rm Flg}(\ell|w|)^{-4}
                                    e^{-\gamma_\Sigma\ell|w|}
 \le2S_4C_{\rm Flg}\ell^{-1}e^{-\gamma_\Sigma\ell}.
\]
Thus $C_\Sigma=2S_4C_{\rm Flg}$.

Finally constant connections are twists by multiplication
by $e^{\mathrm i\theta\cdot c/\ell_*}$.
Every rectilinear path contributing to the average has
the same constant-connection integral, so its phase is
$e^{\mathrm i\theta\cdot(c-y)/\ell_*}$.
Consequently the covariant chamber averages and fine
inverse are conjugate to $Q,Q^*$ and $T_\theta^{-1}$,
and $D_k(c_\theta)$ is conjugate to $D_\theta^{\rm bl}$.
Dilating the mesh $\tau$ to one gives
\[
 \det(D+\mu\text{ on }\mathcal H_\theta)
          =\tau^{-192n^3}\det(D^{\rm ch}[qH_\theta]+\tau\mu),
                \qquad n=\ell_*\mathfrak N .
\]
Here $\tau\mu_*=\bar m\eta$.
Divide \eqref{eq:term-free-block-determinant} at the two
twists. The factors $b^{4\ell_*^3}$ and
$\tau^{-192n^3}$ cancel, proving
\eqref{eq:term-twist-determinant-identity}.
\end{proof}

\section{The terminal remainder and partition comparison}
\label{subsec:term-remainder}

Use the constant already fixed in the construction
\cite[(25)]{QED3Integration}:
\[
 C_0=\max\{C_P^{\rm ph}c_{\rm bd}(\gamma_P),\,
               C^F,\,4b_++16C^Fc(\gamma^F)\},\qquad b_+=1,
\]
where $c(\gamma)=\sum_{z\in\Z^3}e^{-\gamma|z|_\infty}$ and
$c_{\rm bd}(\gamma)=3e^{3\gamma}c(\gamma)$.
Here $C_P^{\rm ph},\gamma_P$ are the photon kernel constants
and $C^F,\gamma^F$ are the fermion constants. Thus $C_0\ge1$
bounds the photon row sums, both field maps, and the Dirac form
proved in \cref{subsec:interacting-terminal-norms}.
Put $b_-=1-L^{-1}$ and let
$K_{\mathrm{pol}}=e^{\kappa_0}/36$, $\kappa_0=54\log72$,
be the scalar polymer-sum constant of
\cite[Appendix~B, (81)]{NormalizedContraction}. For $\kappa\ge\kappa_0$,
\[
 c_\kappa=\sup_\square\sum_{X\supset\square}
                                      e^{-\kappa d_M(X)}\le K_{\mathrm{pol}}.
\]
The fixed scheme gives the following inequalities at every
allowed stop:
\begin{align}
 8C_0^2\ell_*^6(b_-^{-1}+\mu_*^{-1})
                   &\le e_*^{-\varepsilon/2},\nonumber\\
 4(2C_0+c_\kappa+5c_{\rm fl}+6)\ell_*^3
                   &\le e_*^{-\varepsilon},\nonumber\\
 e_*^{1/4-9\varepsilon}&\le\tfrac18,\qquad
 4C_0\ell_*^2p_*<e_*^{-\varepsilon},\nonumber\\
 12\ell_*^3\left(e^{-p_{0,*}/2}
                         +e^{-p_*^2/(2C_0^2)}\right)&\le e_*,
 \qquad
 C_\Sigma e^{-\gamma_\Sigma\ell_*}\le e_*.
                                      \label{eq:term-stop-conditions}
\end{align}
Indeed, the scalar inventory \cite[(29)]{QED3Integration}
holds for every $v\ge-\log e_1$. Write $\ell(v)=LMv^2$.
By \eqref{eq:tpm-stop-scales}, $v_*\ge-\log e_1$,
$\ell_*<\ell(v_*)$, $\mu_*^{-1}<L\ell(v_*)$, and
$\ell_*\ge Mv_*^2$. At $v=v_*$ the covariance condition follows from
\[
 8C_0^2\ell_*^6(b_-^{-1}+\mu_*^{-1})
 \le8C_0^2\ell(v_*)^6(b_-^{-1}+L\ell(v_*))
 \le e_*^{-\varepsilon/2}.
\]
The same upper bound for $\ell_*$ gives the second and fourth
inequalities of \eqref{eq:term-stop-conditions} from the
corresponding polynomial bounds in the inventory; its third
inequality is the inventory's charge bound. For the two tails,
the supplied estimates give
\begin{align*}
12\ell_*^3\big(e^{-p_{0,*}/2}+e^{-p_*^2/(2C_0^2)}\big)
 &\le12\ell(v_*)^3
       \big(e^{-v_*^9/2}+e^{-v_*^{38}/(2C_0^2)}\big)\le e_*,
\\
C_\Sigma e^{-\gamma_\Sigma\ell_*}
 &\le C_\Sigma e^{-\gamma_\Sigma Mv_*^2}\le e_*.
\end{align*}
This verification applies to every $R\in\mathcal T(e)$ above
its mesh threshold. The same inventory at $v_K$ gives
$e_K^{1/4-9\varepsilon}\le1/8$. All these conditions are
uniform in the cutoff.

Put
\begin{equation}\label{eq:term-remainder-errors}
 \varepsilon_1=\tfrac12e_*^{1/4-9\varepsilon},\qquad
 \varepsilon_2=C_\Sigma\ell_*^{-1}e^{-\gamma_\Sigma\ell_*}.
\end{equation}
Together with \eqref{eq:term-complement-errors},
\eqref{eq:term-stop-conditions} gives
$\varepsilon_1\le1/16$, $\varepsilon_2\le e_*$,
and $\varepsilon_4\le\varepsilon_3\le e_*/2$.

\begin{proposition}\label{prop:term-remainder}
For a controlled real pair, every mean-zero real coarse
field $A$ with $\|dA\|_\infty\le p_*$ and every
$\theta\in\R^3$ satisfy
\begin{equation}\label{eq:term-density-remainder}
 Y_*(A+c_\theta)
     =e^{-\varepsilon_k^{\rm v}\ell_*^3}
                R_\theta^{\rm tw}(1+r_*(\theta,A)),
                         \qquad |r_*(\theta,A)|\le\varepsilon_1.
\end{equation}
The remainder is continuous in $(\theta,A)$ on this set.
\end{proposition}

\begin{proof}
The gauge construction in
\cite[(32)--(33)]{MaxwellBlock}
gives a real $\lambda$ such that
$A^g=A-d\lambda$ satisfies $\|A^g\|_\infty\le2\ell_*^2p_*$.
The curvature bound puts both $A+c_\theta$ and
$A^g+c_\theta$ in $\mathrm{Dom}_k$, by
\cite[Theorem~B]{MaxwellBlock};
\cref{thm:term-density} gives equality of their
terminal densities.
Set $\mathscr A^0=H_k^{\rm ph}A^g$.
Summing each of the three photon kernel bounds gives
a norm at most $C_0\|A^g\|_\infty$; hence
\[
 \|\mathscr A^0\|_{\rm each}
                    \le2C_0\ell_*^2p_*<\tfrac12e_*^{-\varepsilon}.
\]
The photon minimizer preserves constants, so the
background is $\mathscr A=\mathscr A^0+c_\theta$.
For $\rho_*=e_*^{-3/4+4\varepsilon}\ge8$ and
$|t|<2\rho_*$, each norm of $t\mathscr A^0$ is
less than $e_*^{-3/4+3\varepsilon}$, the smallest
analytic-domain threshold. Invariance under real
constant shifts gives
\begin{equation}\label{eq:term-terminal-disc}
 \mathscr A^{(t)}=t\mathscr A^0+c_\theta
                  \in\widetilde R_k\subset\widetilde G_k
                         \qquad(|t|<2\rho_*).
\end{equation}

Let $D^{\rm tw}=D_k(c_\theta)$ and
$C_\theta=(D^{\rm tw})^{-1}$.
The free block formula and its twist conjugacy give
$\|C_\theta\|_{\rm op}\le b_-^{-1}+\mu_*^{-1}$.
At the connection $\mathscr A$, after the fermion
field substitution, put
\begin{align}
 V_*={}&-\langle\bar\Psi,
          [D_k(\mathscr A)-D_k(c_\theta)]\Psi\rangle
               -m_k\langle\bar\psi,\psi\rangle\nonumber\\
       &+\sum_XE_k(X;\mathscr A,\psi,\chi).
                                      \label{eq:term-terminal-interaction}
\end{align}
All terms are even. Completing the Grassmann Gaussian
in \eqref{eq:term-Y-density}, at zero external slot, gives
\begin{equation}\label{eq:term-terminal-Gaussian}
 Y_*(A+c_\theta)=e^{-\varepsilon_k^{\rm v}\ell_*^3}
             R_\theta^{\rm tw}\int\dd\mu_{C_\theta}(\Psi)e^{V_*}.
\end{equation}
This is the weight-one identity of the terminal coarse
space; the fine fields retain their $\eta_k^3/48$
pairing in $V_*$.

We estimate the three terms before integrating.
Let $a_1=e_*^{-1/4}$ and
$a_2=C_0^{-1}e_*^{-1/4+\varepsilon}\le a_1$.
The analytic quadratic form
$f(t)=\langle\bar\Psi,D_k(\mathscr A^{(t)})\Psi\rangle$
has $|f(t)|_{a_1}\le C_0e_*^{-1/2}\ell_*^3$ on
\eqref{eq:term-terminal-disc}, by
\eqref{eq:interacting-terminal-Dirac-form}.
Coefficientwise Cauchy's formula with kernel
$[t(t-1)]^{-1}$ on $|t|=\rho_*$ gives
\[
 |f(1)-f(0)|_{a_2}
 \le |f(1)-f(0)|_{a_1}
 \le\frac{C_0e_*^{-1/2}\ell_*^3}{\rho_*-1}
                  \le2C_0e_*^{1/4-4\varepsilon}\ell_*^3 .
\]
The mass-form and activity-substitution bounds
\eqref{eq:interacting-terminal-mass-form} and
\eqref{eq:interacting-terminal-substitution} give
\[
 |m_k\langle\bar\psi,\psi\rangle|_{a_2}
                \le4e_*^{1/4-8\varepsilon}\ell_*^3,\qquad
 |E_k(X;\mathscr A,\psi,\chi)|_{a_2}
                 \le e^{-\kappa d_M(X)}w_k^{\rm E}.
\]
Here $|m_k|\le e_*^{3/4-8\varepsilon}$ and
$w_k^{\rm E}=e_*^{1/4-7\varepsilon}$.
These norm estimates include the ordered pairs within
one chamber cluster. The mass pairing uses the fine site measure
$\eta_k^3/48$;
the four spin components in the mass form therefore
contribute $4\ell_*^3$, rather than $192\ell_*^3$.
Every polymer contains an $M$-cube, so
$\sum_Xe^{-\kappa d_M(X)}
\le(\ell_*/M)^3c_\kappa\le c_\kappa\ell_*^3$.
Consequently
\begin{equation}\label{eq:term-terminal-interaction-bound}
 |V_*|_{a_2}\le
            (2C_0+c_\kappa+4)e_*^{1/4-8\varepsilon}\ell_*^3.
\end{equation}

The covariance integral is the unit-covariance integral
after the one-sided substitution
$\Psi\mapsto C_\theta\Psi$, $\bar\Psi\mapsto\bar\Psi$.
Each covariance row has $4\ell_*^3$ entries of modulus
at most $b_-^{-1}+\mu_*^{-1}$.
Its row norm, together with the identity on barred
variables, is bounded by
\[
 4\ell_*^3(b_-^{-1}+\mu_*^{-1})
 \le\frac{e_*^{-\varepsilon/2}}{2C_0^2\ell_*^3}
 =a_2\,\frac{e_*^{1/4-3\varepsilon/2}}{2C_0\ell_*^3}
 \le a_2.
\]
The weighted substitution estimate therefore sends
$V_*$ to $V_*'$ with $|V_*'|_1\le|V_*|_{a_2}$.
The second inequality of \eqref{eq:term-stop-conditions}
and \eqref{eq:term-terminal-interaction-bound} imply
$|V_*'|_1\le e_*^{1/4-9\varepsilon}/4\le1/32$.
The unit-covariance Wick integral is bounded by the
weight-one Grassmann norm and integrates $1$ to $1$.
Thus
\[
 \left|\int\dd\mu_{C_\theta}e^{V_*}-1\right|
 \le |e^{V_*'}-1|_1
 \le e^{|V_*'|_1}-1
 \le2|V_*'|_1\le\varepsilon_1.
\]
Together with \eqref{eq:term-terminal-Gaussian}, this
proves \eqref{eq:term-density-remainder}.
The density is jointly continuous by
\cref{thm:term-density}. The twist determinant
is a continuous positive function, so its quotient
in \eqref{eq:term-density-remainder} is continuous
on the asserted set.
\end{proof}

\subsection{Partition bounds and comparison of torus sizes}
\label{subsec:term-partition-bounds}

Let $z_{r,n}$ be the matter factor in
\eqref{eq:tpm-matter-factor}, formed at the real mass
$\bar m\eta+m_{\mathrm{in}}$. The fine normalization is
\begin{equation}\label{eq:term-matter-normalization}
 z_{r,n}
 =e^{\varepsilon_{\mathrm{in}}^{\rm v}n^3}
       \det(D^{\rm ch}[0]+\bar m\eta)\,\bar\nu(1)>0.
\end{equation}
Both signs of the total lattice mass are allowed here.
The positive free reference mass $\bar m\eta$ is used
only in the determinant normalization.

\begin{proposition}\label{prop:term-partition-bounds}
For every controlled real pair on $R$, the following
two-sided estimates hold:
\begin{align}
 e^{-\varepsilon_k^{\rm v}\ell_*^3}
 R_\theta^{\rm tw}(1-\varepsilon_1)(1-\varepsilon_4)
 &\le\nu^\theta(\mathfrak G)
 \le e^{-\varepsilon_k^{\rm v}\ell_*^3}
 R_\theta^{\rm tw}(1+\varepsilon_1),
                                   \label{eq:term-small-field-mass}\\
 (1-\varepsilon_3)\bar\nu(1)
 &\le\bar\nu(\mathfrak G)\le\bar\nu(1).
                                   \label{eq:term-full-mass}
\end{align}
In particular
\begin{align}
 \left|\log z_{r,n}-\log z_{r,n}^{\rm free}
       -\varepsilon_{\mathrm{in}}^{\rm v}n^3
                   +\varepsilon_k^{\rm v}\ell_*^3\right|
 &\le2\varepsilon_1+\varepsilon_2
                              +2\varepsilon_3+2\varepsilon_4
 \le2e_*^{1/4-9\varepsilon},\label{eq:term-partition-free-bound}\\
 \left|\log z_{r,n}
       -[p_{\rm ch}^0(\bar m\eta)
                     +\varepsilon_{\mathrm{in}}^{\rm v}]n^3
                    +\varepsilon_k^{\rm v}\ell_*^3\right|
 &\le\varepsilon_5(\bar mR)+2e_*^{1/4-9\varepsilon}.
                                   \label{eq:term-partition-pressure-bound}
\end{align}
\end{proposition}

\begin{proof}
Apply \eqref{eq:term-Borel-observable} at $a=0$ and
$\Phi=1$, then insert \eqref{eq:term-density-remainder}.
The result is
\[
 \nu^\theta(\mathfrak G)=
 e^{-\varepsilon_k^{\rm v}\ell_*^3}R_\theta^{\rm tw}
       \int\chi_*(A)(1+r_*(\theta,A))\,\dd\mathsf m_*(A).
\]
The left side and the prefactor are real, and the
prefactor is positive. The last integral therefore
equals its real part.
The remainder bound and
$1-\varepsilon_4\le\mathsf m_*(\chi_*)\le1$
give \eqref{eq:term-small-field-mass}.
The first inequality of \eqref{eq:term-complement-bound}
gives
$\bar\nu(\mathfrak G^c)\le\varepsilon_3\bar\nu(1)$.
This and positivity prove \eqref{eq:term-full-mass}.

Multiply \eqref{eq:term-small-field-mass} by
$\det(D^{\rm ch}[0]+\bar m\eta)>0$.
The determinant identity
\eqref{eq:term-twist-determinant-identity} and the
penalty twist bound replace this product by the fine
twisted free determinant, with factors between
$e^{-\varepsilon_2}$ and $e^{\varepsilon_2}$.
Integrating in $\theta$ gives
\begin{align*}
 &e^{-\varepsilon_k^{\rm v}\ell_*^3}z_{r,n}^{\rm free}
       e^{-\varepsilon_2}(1-\varepsilon_1)(1-\varepsilon_4)
 \\
 &\hspace{15mm}\le
     \det(D^{\rm ch}[0]+\bar m\eta)\bar\nu(\mathfrak G)
 \le e^{-\varepsilon_k^{\rm v}\ell_*^3}z_{r,n}^{\rm free}
                              e^{\varepsilon_2}(1+\varepsilon_1).
\end{align*}
All integrands are continuous in $\theta$, and the
factors being compared are positive.
Use \eqref{eq:term-full-mass} and
\eqref{eq:term-matter-normalization}, then take logarithms.
Writing the expression on the left of
\eqref{eq:term-partition-free-bound} without absolute
values as $Q_n$, we obtain
\[
 -\varepsilon_2+\log(1-\varepsilon_1)+\log(1-\varepsilon_4)
 \le Q_n\le
 \varepsilon_2+\log(1+\varepsilon_1)-\log(1-\varepsilon_3).
\]
The errors lie in $[0,1/2]$. The inequalities
$-\log(1-x)\le2x$ and $\log(1+x)\le x$ therefore give
the first bound in \eqref{eq:term-partition-free-bound}.
Its right side is at most
$e_*^{1/4-9\varepsilon}+3e_*$.
Since $e_*^{1/4}\le1/8$, we have
$3e_*\le e_*^{1/4}\le e_*^{1/4-9\varepsilon}$,
which gives the second bound.
Adding \eqref{eq:term-free-average-error} proves
\eqref{eq:term-partition-pressure-bound}.
\end{proof}

\subsection{Proof of the common-pressure theorem}
\label{subsec:tpm-common-proof}



\begin{proof}[Proof of \Cref{thm:term-common-partition}]
The common pair is controlled on every torus in the
asserted range, by \cite[Theorem~D]{CoverTuning}.
Its terminal vacuum obeys
$|\varepsilon_k^{\rm v}|
\le5c_{\rm fl}e_*^{1/4-6\varepsilon}$.
The second inequality of \eqref{eq:term-stop-conditions}
therefore gives
\[
 |\varepsilon_k^{\rm v}|\ell_*^3
          \le\tfrac14e_*^{1/4-7\varepsilon}
          \le e_*^{1/4-9\varepsilon}.
\]
Remove this term from
\eqref{eq:term-partition-free-bound} and
\eqref{eq:term-partition-pressure-bound}. This proves
the first two estimates in
\eqref{eq:term-common-pressure-bound}.
The function $\varepsilon_5$ is decreasing.
The scheme's scale bounds give
$\bar mR\ge\mu_K>1/L$ and $e_*\le e_K$.
Since $1/4-9\varepsilon>0$, they prove its last estimate.

Both $R_*$ and $LR_*=R_{\mathrm{top}}$ belong to
$\mathcal T(e)$. The mesh threshold is monotone:
$N_{\min}(LR_*,e)\le N_{\min}(R_*,e)=N_{\rm ref}(e)$.
Thus the common pair and pressure are available on
both tori at every cutoff in the theorem.
The bulk pressure cancels exactly:
\[
 D_{\rm mat}=
 L^3(\log z_{r,n_*}-p_Fn_*^3)
             -(\log z_{r,Ln_*}-p_F(Ln_*)^3).
\]
The pressure bounds give the first inequality in
\eqref{eq:term-matter-comparison}.
Monotonicity of $\varepsilon_5$ and $e_*\le e_K$
give the second. Finally
$\bar mR_*>1/L$ and
$e_K^{1/4-9\varepsilon}\le1/8$ show that the bracket
is at most $\varepsilon_5(1/L)+1$.
The terminal scales, $R_*$ and $e_K$, are independent
of $N$; this proves the stated uniformity.
\end{proof}


\section{Terminal fermion estimates}
\label{subsec:interacting-terminal-operators}

We prove the operator and coefficient estimates used in
the terminal remainder. The reference mass remains variable
throughout this section; the interaction mass is part of the
Grassmann polynomial in \eqref{eq:term-terminal-interaction}.
The block operators and their field domains are those of
\cite[Theorem~A and \S4.2]{FermionBlock} and \cite[\S6]{MaxwellBlock}.

At the stop level write $\eta_*=L^{-k_*}$,
$T_*=\eta_*\Z^3/\ell_*\Z^3$ and
$T_*^0=\Z^3/\ell_*\Z^3$. The fine chamber spinors
have counting measure $\eta_*^3\,48^{-1}\sum_{c,g}$,
and the coarse spinors have four components per unit
site and weight one. We use
\[
 \begin{gathered}
 k_*\ge1,\qquad \ell_*>2M+2,\qquad \gamma^FM\ge16,\\
 1-L^{-1}=:b_-\le b_*\le1,\qquad
 0<\mu_*\le1.
 \end{gathered}
\]
The running mass $\mu_*$ is positive; the bare mass
coupling is treated separately. Let
\begin{align}
 F_*(A)&=D_A+\mu_*+b_*JP_*(A)J^*,&
 S_*(A)&=F_*(A)^{-1},\notag\\
 H_*(A)&=b_*S_*(A)JQ_*^T(-A),&
 D_*(A)&=b_*-b_*Q_*(A)J^*H_*(A).
                                  \label{eq:interacting-terminal-operators}
\end{align}
These are the blocked operators already constructed
in \cite[(51)]{FermionBlock} and
\cite[(52)]{FermionBlock}, at $k=k_*$.
The symbol $F_*$ here is an operator on spinors,
distinct from the curvature variable in the external
partition function.

For a unit site $y$ let $\Delta_y$ be its fine block.
The averaging phases have the form
$\varphi_A(y,c)=e^{i\phi_A(y,c)}$, where
$\phi_A=e_*\eta_*\tau_{k_*}A$ is linear in $A$.
The nested block centers give
\begin{equation}\label{eq:interacting-phase-bounds}
 |\phi_B(y,c)|\le\tfrac32e_*|B|_\infty,\qquad
 |\upsilon_B(c,c')|\le3|c-c'|e_*|B|_\infty.
\end{equation}
The first phase uses only bonds inside $\Delta_y$;
the second, which defines $U_A(c,c')=e^{i\upsilon_A(c,c')}$,
uses only bonds in the coordinate box between its endpoints.
For the first bound, the path lengths at block scales
$L^j\eta_*$ add up to
$\frac32(L-1)\eta_*\sum_{j<k_*}L^j\le\frac32$.
Each of the six coordinate-order paths between
$c$ and $c'$ has length at most $3|c-c'|$,
which gives the second bound.

Fix the exponent $1-\varepsilon<\alpha<1$ from the construction.
For distinct chamber sites $x=(c,g)$ and $x'=(c',g')$, the
pair convention of \cite[\S2.2]{NormalizedContraction} is
\[
 d_X(x,x')=\max\{|c-c'|_\infty,\eta_*\}\le M,\qquad
 (\delta_{\alpha,A}v)(x,x')
   =\frac{U_A(c,c')v(x')-v(x)}{d_X(x,x')^\alpha}.
\]
Distances use the shortest simultaneous periodic lift.
For pairs within one cluster, $d_X=\eta_*$ and $U_A=1$.
The barred quotient uses $U_A^{-1}$.
The admissible pairs for each polymer, including its boundary
pairs of chamber bond type $3$, are as in that convention.

For a complex direction $B$, let $P_B$ be the unit sites
$y$ whose block $\Delta_y$ contains an endpoint of a bond
in $\supp B$. Define
\[
 \epsilon_B=e_*|B|_\infty,\qquad
 P_B^+=\{y:d(y,P_B)\le M+1\}.
\]
The distance is the periodic sup distance on the unit
lattice. For an operator $O$ between fine, coarse or
pair spaces, localize each fine spinor in its unit
block and each pair at its first point, and put
\begin{align}
 |O|(y,y')&=\sup_{\|v\|_\infty\le1}
                         \|\mathbf1_yO\mathbf1_{y'}v\|_\infty,\notag\\
 \|O\|_\gamma&=\max\left\{
       \sup_y\sum_{y'}|O|(y,y')e^{\gamma d(y,y')},
       \sup_{y'}\sum_y|O|(y,y')e^{\gamma d(y,y')}\right\}.
                                  \label{eq:interacting-operator-majorants}
\end{align}
The spinor sup norm uses the maximum of component
moduli. Inserting the unit-block partition between
two factors gives
$|O_1O_2|\le|O_1||O_2|$ and
$\|O_1O_2\|_\gamma\le\|O_1\|_\gamma\|O_2\|_\gamma$.
Write $\|\cdot\|_{\mathrm{op}}$ for the Hilbert-space operator norm
and $\|\cdot\|_{\mathfrak S_1}$ for the trace norm. For coarse
spinor operators,
\begin{equation}\label{eq:interacting-matrix-norms}
 \|O\|_{\mathrm{op}}\le2\|O\|_0,\qquad
 \|O\|_{\mathfrak S_1}\le4\sum_{y,y'}|O|(y,y').
\end{equation}
Indeed the matrix row sums are bounded by $\|O\|_0$
and the column sums by $4\|O\|_0$, so Schur's test
gives the first bound. Polar decomposition bounds
the trace norm by the sum of all entry moduli, and
the entries in one block have total modulus at most
four times its maximum row sum.

We use the fixed constants $C^F,\gamma^F$
of the fermion construction, with its calibration
$a'\ge1$ and
\[
 C^F\ge608\exp(1)\max(1,b)a',
 \qquad C^F\ge64\exp(1)b^2a',\qquad b=1.
\]
The first inequality also dominates the component
norm bounds $64a'$ and $608a'$ for the fine
inverse and its short covariant quotients.
The charge ceiling satisfies
$6Le_*^\varepsilon\le1$ and
$3Me_*^\varepsilon\le1$; on $\widetilde G_*$,
$\mathfrak i(A)<e_*^\varepsilon$. The phase
estimates therefore give
$|\varphi_A^{\pm1}|,|U_A^{\pm1}|\le\exp(1)$,
including every allowed long pair.
The unweakened estimate
\cite[(43)]{FermionBlock}, at its
spectral parameter $y=\infty$, and the block
phases give the following terminal bounds.
On the outer field domain $\widetilde G_*$ the inverse
$S_*$ is analytic, and, with
$E(y,y')=e^{-\gamma^Fd(y,y')}$,
\begin{equation}\label{eq:interacting-terminal-input}
 |S_*|,\quad|\delta_{\alpha,A}S_*|,\quad
 |H_*|,\quad|\delta_{\alpha,A}H_*|
                                  \le C^FE.
\end{equation}
The two pair majorants are restricted to $d_X<1$,
including pairs inside a chamber cluster.
The field domain
\[
 \mathcal W=\{A_r+\mathscr C:
          A_r\in\widetilde R_*\text{ real},\
          \mathscr C\in\mathfrak X_*\}
\]
is contained in $\widetilde G_*$, is invariant under
$A\mapsto-A$, and satisfies
$\mathfrak i(A):=e_*|\operatorname{Im}A|_\infty<1/16$
and $3M\mathfrak i(A)<1/4$. For long pairs,
$|\delta_{\alpha,A}v(x,x')|_c
\le\mathrm e|v(x')|_c+|v(x)|_c$.
The constants $C^F\ge1$ and $0<\gamma^F\le1$ are
uniform in the mesh, allowed torus and running mass. These
are the stop-level forms of the walk and local inverse estimates
of \cite[\S3]{FermionBlock}, uniform for $0<\mu_*\le1$. To obtain the bounds
for $H_*$, write
$H_*=b_*S_*JQ_*^T(-A)$. The last factor has
Euclidean spinor sup norm at most $\exp(1)$,
preserves each unit block, and has no chamber
factor: it repeats the same coarse spinor over
the $48$ labels, multiplied by the block phase.
Composition with the localized inverse bounds
gives constants $32\exp(1)ba'$ and
$304\exp(1)ba'$. Since
$|v|_c\le|v|_2\le2|v|_c$ on $\C^4$,
their component sup norms are at most twice
these constants, hence at most $C^F$.

We shall also need the global component row
bound, without an exponential weight:
\begin{equation}\label{eq:interacting-global-field-rows}
 |H_*^\#(A)|_{1,\infty}\le C^F,\qquad
 |\bar H_*^\#(A)|_{1,\infty}\le C^F .
\end{equation}
Here the rows include every allowed pair, and
the columns are the four coarse components
at the unit sites. Indeed, the unweighted
consequences of \cite[(43)]{FermionBlock}
give
\[
 \|H_*(A)v\|_{\mathrm{su}}
       \le32\exp(1)ba'\|v\|_{\mathrm{su}},\qquad
 \|\delta_{\alpha,A}H_*(A)v\|_{\mathrm{su}}
       \le304\exp(1)ba'\|v\|_{\mathrm{su}}
\]
for short pairs. These include distinct
labels at the same photon site.
For a long pair, $1\le d_X(x,x')\le M$,
so $d_X^{-\alpha}\le1$ and $|U_A|\le\exp(1)$.
Consequently
\[
 |(\delta_{\alpha,A}H_*(A)v)(x,x')|_2
 \le(1+\exp(1))\|H_*(A)v\|_{\mathrm{su}}
 \le128\exp(1)ba'\|v\|_{\mathrm{su}} .
\]
Thus $304\exp(1)ba'$ bounds the Euclidean
sup norm for all pair rows. Conversion to
the component sup norm costs a factor of
two, and the calibration of $C^F$ gives
\eqref{eq:interacting-global-field-rows}.
To identify this operator norm with the row
sum, fix one output component and prescribe
each input component to have the conjugate
phase of its matrix entry. The coarse
spinor values on a fundamental domain can
be prescribed independently; their
antiperiodic extension does not affect
the entry moduli. All unit sites have
weight one.

The barred bound follows from
$\bar H_*(A)=\mathsf C^{-1}H_*(-A)\mathsf C$.
The charge-conjugation matrix $\mathsf C$
is a signed permutation of the four
components, and commutes with the scalar
phases in the covariant quotient. It
preserves both spinor norms. Applying
the preceding argument at $-A$ proves
the second inequality in
\eqref{eq:interacting-global-field-rows}.
Finally, localization of the long-pair
argument gives the bound
$C^F\exp(-\gamma^F\min\{d([x],y'),d([x'],y')\})$
for the coefficient at a coarse site $y'$.
The intermediate Euclidean constant is
$128\exp(1)ba'$, and its component constant
$256\exp(1)ba'$ is again at most $C^F$.

\subsection{The real covariance}
\label{subsec:interacting-covariance}

The covariance at a real field requires no
small-field hypothesis. If $A_r$ is real,
$R=D_{A_r}+\mu_*$ has real part at least $\mu_*$.
The averaging map $Q=Q_*(A_r)J^*$ satisfies
$Q^*=JQ_*^T(-A_r)$ and $QQ^*=1$, because its
phases have modulus one and every block has
fine measure one. Thus $F_*(A_r)=R+b_*Q^*Q$
also has real part at least $\mu_*$.
Set $C=QR^{-1}Q^*$ and $Y=(b_*^{-1}+C)^{-1}$.
The inverse $R^{-1}$ is accretive, so
$b_*^{-1}+C$ is invertible. Multiplication
verifies
\[
 S_*=R^{-1}-R^{-1}Q^*YQR^{-1},\qquad
 Q S_*Q^*=b_*^{-1}(1-b_*^{-1}Y).
\]
The definition of $D_*$ therefore gives
\begin{equation}\label{eq:interacting-coarse-covariance}
 D_*(A_r)^{-1}=b_*^{-1}+QR^{-1}Q^*,\qquad
 \|D_*(A_r)^{-1}\|_{\mathrm{op}}\le b_*^{-1}+\mu_*^{-1}.
\end{equation}

Its block decay follows by conjugating the fine
inverse with a weight. Fix $y'$ and put
$\psi(c)=\vartheta\,d(c,\Delta_{y'})$ on the
periodic fine lattice, with
$\vartheta=\mu_*/(12\mathrm e)\le1$.
This weight is zero on $\Delta_{y'}$, is
$\vartheta$-Lipschitz, and on $\Delta_y$ is
at least $\vartheta(d(y,y')-1)$.
Only the external Dirac hop changes under
conjugation by $e^\psi$, since the internal
hops stay over one photon site. Consequently
\[
 \|e^\psi D_{A_r}e^{-\psi}-D_{A_r}\|_{\mathrm{op}}
 \le6\eta_*^{-1}(e^{\vartheta\eta_*}-1)
 \le6\mathrm e\vartheta=\mu_*/2.
\]
The conjugated $R$ has real part at least
$\mu_*/2$, so its inverse has norm at most
$2/\mu_*$. Apply it between $Q^*$ times
unit spinors at $y$ and $y'$; these have
norm one and supports in their respective
blocks. Moving the weight across the
scalar product gives
\begin{equation}\label{eq:interacting-real-covariance-decay}
 \|D_*(A_r)^{-1}(y,y')\|_{\mathrm{op}}
  \le b_*^{-1}1_{\{y=y'\}}+
       \frac{2\mathrm e}{\mu_*}
                       e^{-\mu_*d(y,y')/(12\mathrm e)}.
\end{equation}
The factor $\mathrm e$ accounts for the unit
block diameter in the lower bound for $\psi$.
\subsection{The terminal Grassmann norms}
\label{subsec:interacting-terminal-norms}

The terminal integral is over four coarse
components per unit site, with Grassmann
weight one. For a polynomial
$F=\sum_I F_I\Psi_I$ in ordered distinct generators,
write $|F|_a=\sum_I|F_I|a^{|I|}$.
The fine norm $\|F\|_{h_*}$ uses the prescribed fine measure
and the two weights for spinor and H\"older-pair generators,
as in \cite[(4)--(5)]{NormalizedContraction};
the weighted substitution estimate is proved in
\cite[Proposition~5.1]{FermionBlock}.
Set
\[
 a_1=e_*^{-1/4},\qquad
 a_2=C_0^{-1}e_*^{-1/4+\varepsilon},\qquad
 h_*=(e_*^{-1/4},e_*^{-1/4+\varepsilon}).
\]
By \eqref{eq:interacting-global-field-rows},
the substitution $H_*^\#(A)\Psi$ has row
norm at most $C^F\le C_0$, for both spinor
and H\"older legs, including all long
and cluster pairs. The coarse kernel
$b_*-D_*(A)=b_*^2Q_*J^*S_*JQ_*^T(-A)$
is bounded by $C^FE(y,y')$: block
compression gives Euclidean constant
$32\exp(1)b^2a'$ by
\cite[(60)]{FermionBlock}, and
conversion to the component norm doubles
it. The calibration of $C^F$ applies.

The weighted Grassmann substitution estimate
applies with $C_0a_2=h_{*,2}\le h_{*,1}$.
Therefore
\begin{equation}\label{eq:interacting-terminal-substitution}
 |F(H_*^\#(A)\Psi)|_{a_2}\le\|F\|_{h_*}.
\end{equation}
For the mass form before substitution, the
fine measure sums to $\ell_*^3$ and there
are four spin components. Its norm at
weight $C_0a_2$ is exactly
$4\ell_*^3(C_0a_2)^2$. In particular
\begin{equation}\label{eq:interacting-terminal-mass-form}
 |\langle\bar\psi,\psi\rangle|_{a_2}
                   \le4e_*^{-1/2+2\varepsilon}\ell_*^3.
\end{equation}
The $48$ labels are canceled by the fine
measure $1/48$; there is no factor $192$
in this estimate. Similarly the coarse
bilinear form has norm $a_1^2$ times the
sum of its matrix entry moduli. Its diagonal
term contributes $4b_*$ per unit site,
and its decaying kernel contributes at
most $4C^F\sum_{y'}E(y,y')$.
For $c(\gamma)=\sum_{z\in\Z^3}e^{-\gamma|z|_\infty}$,
the lattice-sum estimate in \cref{sec:tpm-marked} gives
$c(\gamma)\le(12/\gamma)^3$ when $0<\gamma\le6$.
The chosen $C_0\ge4+16C^Fc(\gamma^F)$ therefore gives
\begin{equation}\label{eq:interacting-terminal-Dirac-form}
 |\langle\bar\Psi,D_*(A)\Psi\rangle|_{a_1}
                                  \le C_0e_*^{-1/2}\ell_*^3.
\end{equation}
These estimates hold on $\widetilde G_*$.
For the external activity evaluated at
$\mathscr T=\mathscr A+\mathscr C$, the
same substitution applies because
$\widetilde R_*+\mathfrak X_*\subset\widetilde G_*$.


\appendix
\section{Marked insertions and reference-mass covariance}
\label{sec:tpm-marked}

The terminal operators of \cref{subsec:interacting-terminal-operators}
also give estimates for local variations of the connection.
These estimates are used for the differentiated terminal integral.
The first subsection uses the field domain $\mathcal W$ and
the operator bounds already stated. For the neighborhood and
covariance constants in the second subsection we additionally fix
\[
                    L^{-3}=:\mu_-<\mu_*\le L^{-1}.
\]
This interval holds at the reference stop; it is not imposed
in \cref{thm:term-common-partition}.

\subsection{Local derivatives of the block operator}

Put
\[
 \mathcal S=C^F(24/\gamma^F)^3,\qquad
 h_{\mathrm L}=6\mathcal S,\qquad
 C_H=203M(C^F)^2,\qquad C_K=226\mathcal S^2.
\]
The lattice sum
$\sum_{y'}e^{-\gamma d(y,y')}\le(12/\gamma)^3$,
$0<\gamma\le6$, follows by choosing nearest periodic
representatives and summing
$e^{-\gamma|z|_\infty}\le e^{-\gamma|z|_1/3}$ over
$\Z^3$. The one-dimensional sum is
$\coth(\gamma/6)\le12/\gamma$.

For a fine spinor or Hölder-pair index $u$, let
$\operatorname{Leg}_A(u;y',\beta')$ be the coefficient
of $\Psi_{\beta'}(y')$ after substituting
$H_*^\#(A)\Psi$ in that generator. Barred generators
use the corresponding barred coefficient.
The set $\operatorname{pos}(u)$ consists of the unit
block containing the first endpoint, and, for a
pair with $d_X(x,x')\ge1$, also the block
containing its second endpoint. Thus it has at
most two elements.

\begin{proposition}\label{prop:interacting-local-insertions}
For $A\in\mathcal W$, the blocked field maps and
their Hölder maps are analytic in $A$. The first
coarse-operator derivative satisfies
\begin{align}
 |\partial_BD_*(A)|(y,y')
  &\le226(C^F)^2\epsilon_B
       \sum_{z\in P_B}E(y,z)E(z,y'),\notag\\
 \|\partial_BD_*(A)\|_{\gamma^F/2}
  &\le C_K\epsilon_B,\qquad
 \sum_{y,y'}|\partial_BD_*(A)|(y,y')
                           \le C_K\epsilon_B|P_B|.
                                  \label{eq:interacting-coarse-first}
\end{align}
If $d(P_B,P_{B'})\ge M+2$, then
\begin{align}
 |\partial_B\partial_{B'}D_*(A)|(y,y')
 \le{}& C_H^2\epsilon_B\epsilon_{B'}
       \sum_{z\in P_B,z'\in P_{B'}}
       \big(E(y,z)E(z,z')E(z',y')\notag\\*
       &\hspace{38mm}+E(y,z')E(z',z)E(z,y')\big).
                                  \label{eq:interacting-coarse-mixed}
\end{align}
The corresponding fine-field legs satisfy
\begin{align}
 |\operatorname{Leg}_A|(u;y')
   &\le3C^F\sum_{w\in\operatorname{pos}(u)}E(w,y'),\notag\\
 |\partial_B\operatorname{Leg}_A|(u;y')
   &\le C_H\epsilon_B
       \sum_{w\in\operatorname{pos}(u),z\in P_B^+}E(w,z)E(z,y'),
                                  \label{eq:interacting-leg-first}\\
 |\partial_B\partial_{B'}\operatorname{Leg}_A|(u;y')
   &\le C_H^2\epsilon_B\epsilon_{B'}
       \sum_{\substack{w\in\operatorname{pos}(u)\\z\in P_B^+,z'\in P_{B'}^+}}
       \big(E(w,z)E(z,z')E(z',y')\notag\\
   &\hspace{37mm}+E(w,z')E(z',z)E(z,y')\big).
                                  \label{eq:interacting-leg-mixed}
\end{align}
The last bound has the same separation hypothesis.
These estimates hold for barred and unbarred legs.
\end{proposition}

\begin{proof}
The block averages and their transposes are diagonal
in unit blocks. From
\eqref{eq:interacting-phase-bounds} and
$\mathfrak i(A)\le1/4$ their majorants are at most
$3/2$. Their directional derivatives are bounded by
\[
                 (9/4)\epsilon_B1_{\{y\in P_B\}}.
\]
Differentiating the only phased Dirac hop gives
\[
 (\partial_BD_Av)(x)=
        -ie_*B(c,c_3)M_3(x)u_A(c,c_3)v(x^{(3)}).
\]
Its majorant is at most $8\epsilon_B$, vanishes
unless the unit blocks coincide or are neighbors,
and vanishes unless both belong to $P_B$.
The derivative of $JP_*(A)J^*$ has diagonal
majorant at most $7\epsilon_B$. Thus for
$F_B=\partial_BF_*(A)$,
\[
 |F_B|(z,z')\le\epsilon_B
       (15\,1_{\{z=z'\}}+8\,1_{\{z,z'\text{ neighbors}\}})
                                      1_{\{z,z'\in P_B\}}.
\]
Summing the six neighboring blocks gives
$\|F_B\|_\gamma\le146\epsilon_B$ for $0\le\gamma\le1$.
Also, since $E(z',y')\le\mathrm eE(z,y')$ for
neighboring $z,z'$,
\begin{equation}\label{eq:interacting-insertion-chain}
 \sum_{z,z'}E(y,z)|F_B|(z,z')E(z',y')
        \le146\epsilon_B\sum_{z\in P_B}E(y,z)E(z,y').
\end{equation}

If the supports are separated by $M+2$, the mixed
derivatives of $D_A$, the averages, $JP_*J^*$,
$F_*$ and $\delta_{\alpha,A}$ vanish. Each entry
is an exponential of a linear phase supported
in a unit block, a single external bond or a
coordinate box of diameter at most $M$.
Two separated directions cannot both occur in
one phase. For $JP_*J^*$ the remaining product
terms have disjoint unit-block supports.

The inverse derivative and the blocked-field
definition give
\[
 \partial_BH_*=-S_*F_BH_*+
               b_*S_*\partial_B(JQ_*^T(-A)).
\]
Use \eqref{eq:interacting-terminal-input} and
\eqref{eq:interacting-insertion-chain} in the first
term. The diagonal derivative in the second
is absorbed as the summand $z=y'$. The combined
coefficient is at most $149(C^F)^2$:
\begin{equation}\label{eq:interacting-H-first}
 |\partial_BH_*|(y,y')
 \le149(C^F)^2\epsilon_B
                    \sum_{z\in P_B}E(y,z)E(z,y').
\end{equation}
The same estimate holds for
$\delta_{\alpha,A}\partial_BH_*$ on short pairs,
since the first factor $S_*$ can be replaced
by $\delta_{\alpha,A}S_*$.

For separated directions differentiating once more
leaves four terms:
\begin{align*}
 \partial_B\partial_{B'}H_*={}&
 S_*F_BS_*F_{B'}H_*+S_*F_{B'}S_*F_BH_*\\*
 &-b_*S_*F_BS_*\partial_{B'}(JQ_*^T(-A))
  -b_*S_*F_{B'}S_*\partial_B(JQ_*^T(-A)).
\end{align*}
Apply \eqref{eq:interacting-insertion-chain} twice
to the first two terms. Each gives an ordered
three-link exponential chain with coefficient
$146^2(C^F)^3$. The diagonal last factors in
the other two terms give the same chains, with
coefficient $(657/2)(C^F)^2$. Each ordered
chain therefore has coefficient at most
$21645(C^F)^3$. Again the first factor may
be replaced by the short-pair Hölder inverse.

Differentiate $D_*=b_*-b_*Q_*J^*H_*$.
The derivative of the average contributes
$(9/4)C^F$ to the first exponential chain and
the derivative of $H_*$ contributes
$(3/2)149(C^F)^2$. Their sum is at most
$226(C^F)^2$, proving the first line of
\eqref{eq:interacting-coarse-first}.
Weighted row and column sums give two factors
$(24/\gamma^F)^3$, hence $C_K$; without weights
the same sum retains one factor $|P_B|$.
For the mixed derivative the three product
terms contain one derivative on each factor,
or the mixed derivative of $H_*$. Each chain
has coefficient at most
$(1341/4)(C^F)^2+(64935/2)(C^F)^3
\le32803(C^F)^3\le C_H^2$.
This proves \eqref{eq:interacting-coarse-mixed}.

A leg is a component row of $H_*$ or of
$\delta_{\alpha,A}H_*$. For a short pair its
position is its first unit block; for a long
pair its position set contains both endpoint
blocks. The sum over spin components of the
row is bounded by the relevant majorant.
For long pairs use
$|\delta_{\alpha,A}v|\le\mathrm e|v(x')|+|v(x)|$;
for short pairs use the Hölder input directly.
This gives the first line of
\eqref{eq:interacting-leg-first}. In particular, with
$d(y',\operatorname{pos}(u))=\min_{w\in\operatorname{pos}(u)}d(y',w)$,
\[
 \sum_{y',\beta'}|\operatorname{Leg}_A(u;y',\beta')|
       e^{(\gamma^F/2)d(y',\operatorname{pos}(u))}
 \le3C^F|\operatorname{pos}(u)|(24/\gamma^F)^3
 \le h_{\mathrm L}.
\]
Indeed, multiplying each endpoint decay by this weight leaves
decay rate at least $\gamma^F/2$, and there are at most two endpoints.

The derivative of a Hölder leg has two terms:
$(\partial_B\delta_{\alpha,A})H_*+
\delta_{\alpha,A}\partial_BH_*$.
The first is bounded by
$4d_X^{1-\alpha}\epsilon_B|H_*v(x')|$,
and vanishes unless the pair is within $M$
of $P_B$. It vanishes for pairs inside a
cluster, since their phase is one.
For short pairs this adds at most
$11C^F$ to $149(C^F)^2$; for long pairs
it adds $4MC^F$ to $(4/3)149(C^F)^2$.
Both are at most $C_H$, after enlarging
$P_B$ to $P_B^+$. This proves the derivative
bound in \eqref{eq:interacting-leg-first}.

For its mixed derivative the second derivative
of the Hölder phase vanishes by separation.
There remain
$(\partial_B\delta_{\alpha,A})\partial_{B'}H_*$
and its interchange, and
$\delta_{\alpha,A}\partial_B\partial_{B'}H_*$.
On short pairs each ordered chain has
coefficient at most
$1639(C^F)^2+21645(C^F)^3$.
On long pairs the coefficient at the second
endpoint is at most
$596M(C^F)^2+28860(C^F)^3$;
the first endpoint has the smaller coefficient
$21645(C^F)^3$. All are at most $C_H^2$,
proving \eqref{eq:interacting-leg-mixed}.
Barred legs follow by the monomial charge
conjugation matrix and the replacement $A\mapsto-A$,
which preserve component row sums; the domain
is invariant under that replacement.
\end{proof}


\subsection{A complex neighborhood and covariance bounds}

Choose
\[
 c_\tau=\min\{1/8,\mu_-/48,(292\mathcal S)^{-1}\},\qquad
 C_K'=1328\mathcal S^2,\qquad
 \mathcal W^+=\{A+A_1:A\in\mathcal W,\ e_*|A_1|_\infty\le c_\tau\}.
\]
Every point of $\mathcal W^+$ has imaginary part
$\mathfrak i<1/4$. The elementary averages and
$F_*$ therefore satisfy the same local derivative
bounds as above, although the exponential propagator
input is initially given only on $\mathcal W$.
The enlargement preserves invertibility:
\begin{equation}\label{eq:interacting-enlarged-inverse}
 \|S_*(A+A_1)\|_{\gamma^F/2}\le2\mathcal S.
\end{equation}
To prove this, integrate $\partial_{A_1}F_*$
along the segment to get
$\|F_*(A+A_1)-F_*(A)\|_{\gamma^F/2}
\le146c_\tau$. Multiplying by $S_*(A)$ gives
norm at most $146c_\tau\mathcal S\le1/2$.
The Neumann inverse in the submultiplicative
majorant norm proves \eqref{eq:interacting-enlarged-inverse}.
Since its finite-dimensional determinant is
nonzero, $S_*$ is analytic in a neighborhood
of each point of $\mathcal W^+$, as is $D_*$.

Differentiate
$D_*=b_*-b_*^2Q_*J^*S_*JQ_*^T(-A)$.
The three terms have majorant-norm coefficients
\[
 \tfrac94(2\mathcal S)\tfrac32,\qquad
 \tfrac32(2\mathcal S)146(2\mathcal S)\tfrac32,\qquad
 \tfrac32(2\mathcal S)\tfrac94.
\]
Their sum is
$\frac{27}{2}\mathcal S+1314\mathcal S^2\le C_K'$.
Keeping the support indicator of either the
average derivative or $F_B$ gives, respectively,
\begin{equation}\label{eq:interacting-enlarged-derivative}
 \|\partial_BD_*(A)\|_{\gamma^F/2}\le C_K'\epsilon_B,\qquad
 \sum_{y,y'}|\partial_BD_*(A)|(y,y')
                                  \le C_K'\epsilon_B|P_B|,
                                      \qquad A\in\mathcal W^+.
\end{equation}
In the middle term, sum first the two outer
$S_*$ factors using their row and column bounds
$2\mathcal S$, and then sum the local derivative
over its first endpoint in $P_B$. Integrating
\eqref{eq:interacting-enlarged-derivative} along
any segment inside $\mathcal W^+$ gives the
same two bounds for $D_*(A+B)-D_*(A)$.


Set
\[
 g=(b_-^{-1}+\mu_-^{-1})^{1/2},\qquad
 h_G=1+g,\qquad g_1=b_-^{-1}+4/\mu_-,
 \qquad \theta_{\mathrm{CT}}=\mu_-/(12\mathrm e),
\]
\[
 \vartheta_0=\min\{\theta_{\mathrm{CT}}/2,\gamma^F/8,1/M\},
 \qquad
 \alpha_G=2b_-^{-1}
        +4\mathrm e\mu_-^{-1}(24/\theta_{\mathrm{CT}})^3.
\]
A component row or column of a $4\times4$ matrix
has sum at most twice its Hilbert-space norm.
Using \eqref{eq:interacting-real-covariance-decay}
and the lattice exponential sum, every row and
column of $G=D_*(A_r)^{-1}$ therefore satisfies
\begin{equation}\label{eq:interacting-Gaussian-covariance}
 \sum_{u'}|G(u,u')|e^{\vartheta_0d(u,u')}
       \le\alpha_G,\qquad
 \sum_u|G(u,u')|e^{\vartheta_0d(u,u')}
       \le\alpha_G,\qquad \|G\|_{\mathrm{op}}\le g^2.
\end{equation}
These constants depend on $L,C^F,\gamma^F,M$,
not on the size of the real connection.

The same inverse identity gives a bound for the
coarse covariance at complex fields. Suppose
$A=A_r+A_1$, with $A_r$ real,
$e_*|A_1|_\infty\le\mu_*/24$, and $F_*(A)$
invertible. The external-hop phase difference
gives $\|D_A-D_{A_r}\|_{\mathrm{op}}\le12e_*|A_1|_\infty
\le\mu_*/2$. Hence $D_A+\mu_*$ is accretive
with inverse norm at most $2/\mu_*$. Its two
average maps have norm at most $\sqrt2$,
because $|\varphi_A^{\pm1}|\le e^{1/16}\le\sqrt2$.
The same finite-rank inversion calculation
then gives
\begin{equation}\label{eq:interacting-complex-covariance}
 D_*(A)^{-1}=b_*^{-1}+
 Q_*(A)J^*(D_A+\mu_*)^{-1}JQ_*^T(-A),
 \qquad
 \|D_*(A)^{-1}\|_{\mathrm{op}}\le g_1.
\end{equation}
Here invertibility of $b_*^{-1}+QR^{-1}Q'$
follows from that of $F_*(A)$ and $R$:
$\det(1+b_*R^{-1}Q'Q)=\det(1+b_*QR^{-1}Q')$.
This determinant identity follows by the two
block-triangular factorizations of
$\left(\begin{smallmatrix}1&-b_*R^{-1}Q'\\Q&1\end{smallmatrix}\right)$.
The inverse formula is then checked by
multiplication exactly as above.
Each covariance entry is bounded by $g_1$,
so each row has sum at most $4\ell_*^3g_1$.


\printbibliography
\end{document}
